% ============================================================
% CHƯƠNG 4 — CÔNG NGHỆ, XÂY DỰNG VÀ KIỂM THỬ HỆ THỐNG
% Tạo ngày 17/11/2026 khi Chương 2 đổi thành "Phân tích hệ thống" và Chương 3 thành "Thiết kế hệ thống".
% Chương này gom (i) khung chương "Công nghệ sử dụng" cũ (Chương 2 trước 17/11/2026) và (ii) các mục
% Xây dựng ứng dụng, Kiểm thử, Triển khai của Chương 3 cũ "Thiết kế, triển khai và đánh giá".
% Bản trước: backup/2026-11-17_truoc-phantich-thietke/Chap2.tex, Chap3.tex.
% CHƯA \input trong main.tex (chưa viết nội dung) — bỏ dấu % ở dòng % ============================================================
% CHƯƠNG 4 — CÔNG NGHỆ, XÂY DỰNG VÀ KIỂM THỬ HỆ THỐNG
% Tạo ngày 17/11/2026 khi Chương 2 đổi thành "Phân tích hệ thống" và Chương 3 thành "Thiết kế hệ thống".
% Chương này gom (i) khung chương "Công nghệ sử dụng" cũ (Chương 2 trước 17/11/2026) và (ii) các mục
% Xây dựng ứng dụng, Kiểm thử, Triển khai của Chương 3 cũ "Thiết kế, triển khai và đánh giá".
% Bản trước: backup/2026-11-17_truoc-phantich-thietke/Chap2.tex, Chap3.tex.
% CHƯA \input trong main.tex (chưa viết nội dung) — bỏ dấu % ở dòng % ============================================================
% CHƯƠNG 4 — CÔNG NGHỆ, XÂY DỰNG VÀ KIỂM THỬ HỆ THỐNG
% Tạo ngày 17/11/2026 khi Chương 2 đổi thành "Phân tích hệ thống" và Chương 3 thành "Thiết kế hệ thống".
% Chương này gom (i) khung chương "Công nghệ sử dụng" cũ (Chương 2 trước 17/11/2026) và (ii) các mục
% Xây dựng ứng dụng, Kiểm thử, Triển khai của Chương 3 cũ "Thiết kế, triển khai và đánh giá".
% Bản trước: backup/2026-11-17_truoc-phantich-thietke/Chap2.tex, Chap3.tex.
% CHƯA \input trong main.tex (chưa viết nội dung) — bỏ dấu % ở dòng % ============================================================
% CHƯƠNG 4 — CÔNG NGHỆ, XÂY DỰNG VÀ KIỂM THỬ HỆ THỐNG
% Tạo ngày 17/11/2026 khi Chương 2 đổi thành "Phân tích hệ thống" và Chương 3 thành "Thiết kế hệ thống".
% Chương này gom (i) khung chương "Công nghệ sử dụng" cũ (Chương 2 trước 17/11/2026) và (ii) các mục
% Xây dựng ứng dụng, Kiểm thử, Triển khai của Chương 3 cũ "Thiết kế, triển khai và đánh giá".
% Bản trước: backup/2026-11-17_truoc-phantich-thietke/Chap2.tex, Chap3.tex.
% CHƯA \input trong main.tex (chưa viết nội dung) — bỏ dấu % ở dòng \input{Chap4} để khôi phục.
% Bảng API (\BangAPI) đã dùng ở Chương 3 (Thiết kế API), không đặt lại ở đây.
%
% Hướng dẫn của mẫu SoICT cho phần công nghệ: không quá 10 trang; chỉ giới thiệu công nghệ đồ án thực sự
% dùng; với mỗi công nghệ nêu (i) yêu cầu nó phục vụ (NC-xx, VĐ-x, yêu cầu phi chức năng ở Chương 2),
% (ii) phương án thay thế, (iii) lý do chọn.
% Văn phong: xem ghi chú đầu Chap1.tex. [VIẾT SAU] toàn chương.
% ============================================================

\chapter{CÔNG NGHỆ, XÂY DỰNG VÀ KIỂM THỬ HỆ THỐNG}
\chaptermark{CÔNG NGHỆ, XÂY DỰNG VÀ KIỂM THỬ}
\label{chap:congnghe}
\label{chap:cosolythuyet}
\label{chap:xaydung}
\label{chap:ketqua}
\input{figures_chap4}
\input{figures_chap5}

% Tổng quan (Normal, không tiêu đề): chương trình bày các công nghệ đồ án dùng để hiện thực thiết kế ở
% Chương 3, kết quả cài đặt, kiểm thử và triển khai trên Sepolia.
% [VIẾT SAU] \canviet{Đoạn tổng quan chương}

% ============ PHẦN A — CÔNG NGHỆ SỬ DỤNG (từ Chương 2 cũ) ============
% [VIẾT SAU] Toàn bộ các mục dưới đây chưa viết nội dung — bỏ dấu % để khôi phục.
% % ------------------------------------------------------------
% \section{Blockchain và nền tảng Ethereum}
% \label{sec:blockchain}
% \label{sec:ethereum}
%
% \subsection*{Blockchain}
% \label{subsec:cautrucblock}
% \label{subsec:dongthuan}
% % Ý cần có (ngắn gọn, chỉ phần đồ án dùng): khối, mã băm khối trước, sổ cái chỉ ghi thêm (ledger) ->
% % vì sao dữ liệu on-chain khó sửa, phục vụ yêu cầu toàn vẹn ở mục \ref{sec:phichucnang} và VĐ-2;
% % đồng thuận PoW/PoS, Ethereum chuyển sang PoS (The Merge 15/9/2022).
% % Trích: \cite{nakamoto2008bitcoin,zheng2017overview,ethereum2024pos}. Nhắc \autoref{fig:chuoikhoi} trước hình.
% % [VIẾT SAU] \canviet{Blockchain}
% %
% % \HinhChuoiKhoi
%
% \subsection*{Lựa chọn nền tảng blockchain}
% \label{subsec:phanloai}
% % Bảng so sánh Ethereum (public) với Hyperledger Fabric (permissioned, mà Vaccargiu dùng) và Chia
% % (CAD Trust dùng): mô hình truy cập, đồng thuận, ngôn ngữ smart contract, chuẩn token, ví, testnet, chi phí.
% % Lý do chọn Ethereum: phiếu giao nhiệm vụ; chuẩn ERC-721/ERC-20; MetaMask; testnet công khai để bên thứ ba
% % tự kiểm tra (VĐ-2, NC-07). Nhược điểm: phí gas, dữ liệu công khai -> xử lý ở yêu cầu quyền riêng tư.
% % [VIẾT SAU] \canviet{Lựa chọn nền tảng blockchain}
%
% \subsection*{Tài khoản, giao dịch và phí gas}
% \label{subsec:taikhoan}
% \label{subsec:evm}
% % Ý cần có: EOA và contract account; khóa riêng, địa chỉ ví; giao dịch ký số, nonce; EVM, gas, EIP-1559.
% % Đây là cơ sở cho phần đo gas ở Chương 3. Trích: \cite{buterin2014ethereum,wood2014yellow,eip1559}.
% % [VIẾT SAU] \canviet{Tài khoản, giao dịch và phí gas}
%
% \subsection*{Testnet Sepolia và Etherscan}
% \label{subsec:sepolia}
% % Ý cần có: mainnet và testnet; Sepolia, faucet; Etherscan để tra giao dịch, xác minh (verify) mã hợp đồng.
% % Lý do: phạm vi nguyên mẫu (mục \ref{sec:muctieu}), không tốn phí thật.
% % [VIẾT SAU] \canviet{Testnet Sepolia}
%
% % ------------------------------------------------------------
% \section{Smart contract và công cụ phát triển}
% \label{sec:smartcontract}
%
% \subsection*{Smart contract và ngôn ngữ Solidity}
% \label{subsec:khainiemsc}
% \label{subsec:solidity}
% % Ý cần có: định nghĩa (Szabo); tự thực thi, không sửa được mã sau triển khai; event và log (dùng cho
% % bộ đồng bộ sự kiện); cấu trúc contract (state variable, function, modifier, event), kế thừa.
% % Yêu cầu phục vụ: NC-03, NC-06 (ràng buộc sê-ri, không hủy hai lần thành ràng buộc kỹ thuật).
% % Trích: \cite{szabo1997formalizing,buterin2014ethereum,solidity2024docs}.
% % [VIẾT SAU] \canviet{Smart contract và Solidity}
%
% \subsection*{Thư viện OpenZeppelin và Hardhat}
% \label{subsec:openzeppelin}
% % OpenZeppelin 5: ERC721, ERC20, AccessControl (VERIFIER_ROLE...), Pausable, ReentrancyGuard -> NC-09 và
% % yêu cầu bảo mật. Hardhat: biên dịch, Hardhat Network, test Mocha/Chai, coverage, gas reporter.
% % Phương án thay thế: Foundry, Truffle (đã ngừng phát triển); lý do chọn Hardhat (JavaScript cùng ngôn ngữ backend).
% % Trích: \cite{openzeppelin2024contracts}.
% % [VIẾT SAU] \canviet{OpenZeppelin và Hardhat}
%
% \subsection*{Các lỗ hổng thường gặp của smart contract}
% \label{subsec:baomat}
% % Reentrancy, kiểm soát truy cập sai, tràn số (Solidity >= 0.8 tự kiểm tra), tx.origin. Mỗi lỗ hổng
% % 1–2 câu và cách phòng trong đồ án; công cụ quét Slither. Trích: \cite{atzei2017survey}.
% % [VIẾT SAU] \canviet{Lỗ hổng smart contract}
%
% % ------------------------------------------------------------
% \section{Các chuẩn token}
% \label{sec:token-standards}
%
% \subsection*{ERC-20 và ERC-721}
% \label{subsec:erc20}
% \label{subsec:erc721}
% % ERC-20: fungible token, totalSupply, balanceOf, transfer, approve, transferFrom, allowance \cite{eip20}.
% % ERC-721: NFT, ownerOf, safeTransferFrom, tokenURI, onERC721Received \cite{eip721}.
% % [VIẾT SAU] \canviet{ERC-20 và ERC-721}
%
% \subsection*{Lựa chọn chuẩn token cho tín chỉ các-bon}
% \label{subsec:chontoken}
% % Bảng so sánh ERC-20 / ERC-721 / ERC-1155 theo: giữ metadata riêng từng lô, chia nhỏ, thanh khoản,
% % hỗ trợ ví. Kết luận: mỗi lô là một NFT ERC-721 (giữ dự án, vintage, dãy sê-ri, NC-03), pool ERC-20
% % cho giao dịch từng phần (NC-05); kế thừa mô hình Batch NFT/TCO2/BCT của Toucan (mục \ref{sec:khaosatsobo}).
% % [VIẾT SAU] \canviet{Lựa chọn chuẩn token}
%
% % ------------------------------------------------------------
% \section{Ứng dụng web giao tiếp với blockchain}
% \label{sec:dapp}
%
% \subsection*{Kiến trúc dApp}
% \label{subsec:kientrucdapp}
% % Luồng một giao dịch từ giao diện tới smart contract (giải thích \autoref{fig:luongdapp}).
% % [VIẾT SAU] \canviet{Kiến trúc dApp}
% %
% % \HinhLuongDApp
%
% \subsection*{Ví MetaMask, ethers.js và Next.js}
% \label{subsec:metamask}
% \label{subsec:ethers}
% % MetaMask giữ khóa riêng, EIP-1193 (window.ethereum) \cite{eip1193}; ethers.js: Provider, Signer,
% % Contract, ABI, call và send \cite{ethers2024docs}; Next.js cho giao diện. Yêu cầu phục vụ: tính dễ dùng,
% % người dùng tự ký giao dịch (backend không giữ khóa). Phương án thay thế: web3.js, viem; React thuần.
% % [VIẾT SAU] \canviet{MetaMask, ethers.js và Next.js}
%
% % ------------------------------------------------------------
% \section{Công nghệ phía off-chain}
% \label{sec:offchain-congnghe}
% \label{sec:onoffchain-lythuyet}
% % Node.js + Express (API RESTful), PostgreSQL (dữ liệu theo CAD Trust V2.0, NC-10), JWT, xác thực bằng
% % chữ ký ví. Tiêu chí chọn dữ liệu on-chain/off-chain theo \cite{xu2017taxonomy}; neo mã băm keccak256 của hồ sơ;
% % IPFS \cite{benet2014ipfs} là hướng mở rộng (hiện lưu tệp cục bộ). Phương án thay thế: MongoDB; lý do chọn
% % CSDL quan hệ (mô hình CAD Trust có quan hệ rõ, giống UNDP dùng PostgreSQL).
% % [VIẾT SAU] \canviet{Công nghệ off-chain}


% ============ PHẦN B — XÂY DỰNG, KIỂM THỬ, TRIỂN KHAI (từ Chương 3 cũ) ============
% ------------------------------------------------------------
\section{Xây dựng ứng dụng}
\label{sec:xaydung}

\subsection*{Thư viện và công cụ sử dụng}
\label{sec:moitruong}
% Mẫu SoICT: bảng công cụ có phiên bản (và URL). Giải thích \autoref{tab:congnghe}; lý do chọn công cụ đã nêu
% ở Chương 2 nên chỉ nhắc lại ngắn. Cấu trúc mã nguồn theo \autoref{fig:cautrucma}.
% [VIẾT SAU] \canviet{Thư viện và công cụ sử dụng}

\BangCongNghe

\HinhCauTrucMa

\subsection*{Kết quả đạt được}
\label{subsec:ketquadatduoc}
\label{sec:caidat-hopdong}
\label{sec:caidat-backend}
% Mẫu SoICT: mô tả các sản phẩm đóng gói, thành phần, vai trò; bảng thống kê số dòng mã, số tệp, số hợp đồng,
% số endpoint, dung lượng.
% Nội dung cài đặt từ Chương 4 cũ (chỉ tóm lược, chi tiết đóng góp ở Chương 4):
%  - CarbonCreditNFT: struct CreditBatch (projectId, vintage, serialFrom, quantity, retired, documentHash,
%    verifier, issuedAt); mintBatch: seriesKey = keccak256(projectId, vintage), serialFrom >= nextAvailableSerial,
%    documentUsed[documentHash]; retire hủy từ đầu dãy sê-ri còn lại; retireFor cho pool (RETIREMENT_AGENT_ROLE);
%    Pausable ghi đè _update.
%  - CarbonPool: decimals = 0; nạp qua onERC721Received; redeem đốt remainingOf(tokenId); retire đốt token rồi
%    gọi retireFor; ReentrancyGuard.
%  - Phân quyền: DEFAULT_ADMIN_ROLE và vai trò phát hành ISSUER_ROLE (cơ quan quản lý, mỗi tỉnh một tài khoản),
%    VERIFIER_ROLE (VVB, chỉ phê duyệt off-chain), RETIREMENT_AGENT_ROLE (pool). [CẦN SỬA MÃ] thêm ISSUER_ROLE.
%  - Backend: db/schema.sql từ ERD; xác thực nonce -> personal_sign -> ethers.verifyMessage -> JWT 8 giờ;
%    API theo \autoref{tab:api} (backend không giữ khóa riêng); indexer quét queryFilter, khóa (tx_hash, log_index),
%    chưa xử lý reorg.
% [VIẾT SAU] \canviet{Kết quả đạt được}

% \BangAPI   % [17/11/2026] đã dùng ở Chương 3, mục Thiết kế API

\subsection*{Minh họa các chức năng chính}
\label{subsec:minhhoa}
\label{sec:caidat-frontend}
% Mỗi màn hình 1 ảnh chụp + giải thích ngắn: danh sách dự án, chi tiết dự án, tra cứu sê-ri, chứng nhận hủy,
% dự án của tôi, thẩm định (VVB), ví tín chỉ, dashboard, quản trị. Chụp khi chạy thật (README "Chạy thử cục bộ").
% [VIẾT SAU] \canviet{Minh họa các chức năng chính}

% ------------------------------------------------------------
\section{Kiểm thử}
\label{sec:kiemthu}
\label{sec:pp-kiemthu}
\label{sec:kq-chucnang}
% Mẫu SoICT: 2–3 trang; thiết kế ca kiểm thử cho 2–3 chức năng quan trọng nhất (phát hành, hủy từ pool, tra cứu),
% nêu kỹ thuật kiểm thử; tổng kết số ca và kết quả; giải thích ca không đạt. Ca kiểm thử khác đưa vào Phụ lục.
% Ba mức: kiểm thử đơn vị hợp đồng trên Hardhat Network (Mocha/Chai, loadFixture); kiểm thử API với pg-mem và
% supertest; kiểm thử tích hợp e2e-indexer và smoke-api. Ghi cấu hình máy, phiên bản Node, Hardhat.
% Nhận xét độ phủ: nhánh của CarbonPool 75% vì ... (xem coverage/index.html).
% Chi phí gas: quy đổi ETH/USD theo giá gas tại thời điểm đo (ghi nguồn, ngày); vì sao mintBatch, retire tốn gas.
% Kịch bản chống tính toán kép (\BangTinhToanKep) đặt ở Chương 4, mục \ref{sec:gp-tinhtoankep}.
% [VIẾT SAU] \canviet{Kiểm thử}

\BangTestHopDong

\BangTestAPI

\BangDoPhu

\BangGas

% ------------------------------------------------------------
\section{Triển khai}
\label{sec:trienkhai}
\label{sec:kq-sepolia}
% Giải thích \autoref{fig:trienkhai}: Sepolia + RPC, backend + PostgreSQL, frontend; địa chỉ hợp đồng,
% link Etherscan (đã verify), vài giao dịch mẫu (mint, retire), ảnh chụp.
% [VIẾT SAU] \canviet{Triển khai}

\HinhTrienKhai


% ------------------------------------------------------------
% Kết chương (Normal, không tiêu đề): tóm tắt bộ công nghệ (bảng Công nghệ | Phiên bản | Vai trò | Yêu cầu
% phục vụ), kết quả cài đặt, kiểm thử, triển khai; dẫn sang Chương 5 (giải pháp và đóng góp).
% [VIẾT SAU] \canviet{Kết chương 4}
 để khôi phục.
% Bảng API (\BangAPI) đã dùng ở Chương 3 (Thiết kế API), không đặt lại ở đây.
%
% Hướng dẫn của mẫu SoICT cho phần công nghệ: không quá 10 trang; chỉ giới thiệu công nghệ đồ án thực sự
% dùng; với mỗi công nghệ nêu (i) yêu cầu nó phục vụ (NC-xx, VĐ-x, yêu cầu phi chức năng ở Chương 2),
% (ii) phương án thay thế, (iii) lý do chọn.
% Văn phong: xem ghi chú đầu Chap1.tex. [VIẾT SAU] toàn chương.
% ============================================================

\chapter{CÔNG NGHỆ, XÂY DỰNG VÀ KIỂM THỬ HỆ THỐNG}
\chaptermark{CÔNG NGHỆ, XÂY DỰNG VÀ KIỂM THỬ}
\label{chap:congnghe}
\label{chap:cosolythuyet}
\label{chap:xaydung}
\label{chap:ketqua}
% ============================================================
% HÌNH VÀ BẢNG CHƯƠNG 4 (CÔNG NGHỆ, XÂY DỰNG VÀ KIỂM THỬ) — gọi bằng macro trong Chap4.tex
%   \HinhChuoiKhoi  : cấu trúc liên kết các khối           (fig:chuoikhoi)   [chuyển từ figures_chap2 ngày 17/11/2026]
%   \HinhLuongDApp  : luồng tương tác của dApp              (fig:luongdapp)   [chuyển từ figures_chap2 ngày 17/11/2026]
%   \BangCongNghe   : công nghệ và phiên bản sử dụng      (tab:congnghe)
%   \HinhTrienKhai  : sơ đồ triển khai hệ thống            (fig:trienkhai)
%   \HinhCauTrucMa  : cấu trúc thư mục mã nguồn            (fig:cautrucma)
% Phiên bản lấy từ package.json / hardhat.config.js tại thời điểm 21/10/2026.
% ============================================================

% ---------- Cấu trúc liên kết các khối ----------
\newcommand{\HinhChuoiKhoi}{%
\begin{figure}[htbp]
\centering
\begin{tikzpicture}[
    >=Stealth,
    truong/.style={draw, minimum width=3.3cm, minimum height=0.55cm,
                   font=\footnotesize, fill=white, inner sep=1pt},
    tieude/.style={truong, fill=gray!20, font=\footnotesize\bfseries}]
  \foreach \i/\x/\ten in {0/0/N-1, 1/4.6/N, 2/9.2/N+1} {
    \node[tieude] (h\i) at (\x,0) {Khối \ten};
    \node[truong, below=0pt of h\i] (p\i) {Mã băm khối trước};
    \node[truong, below=0pt of p\i] (t\i) {Dấu thời gian};
    \node[truong, below=0pt of t\i] (m\i) {Gốc Merkle};
    \node[truong, below=0pt of m\i, minimum height=1cm, align=center] (g\i)
        {Danh sách\\giao dịch};
    \node[draw, thick, rounded corners=2pt, inner sep=3pt, fit=(h\i)(g\i)] (b\i) {};
  }
  \draw[->, thick] (p1.west) -- (b0.east |- p1.west);
  \draw[->, thick] (p2.west) -- (b1.east |- p2.west);
  \node[font=\large] at (-2.35,-1.4) {$\cdots$};
  \node[font=\large] at (11.55,-1.4) {$\cdots$};
\end{tikzpicture}
\caption[Cấu trúc liên kết giữa các khối trong chuỗi khối]{Cấu trúc liên kết giữa các khối trong chuỗi khối \hinhtao[Claude]{20/10/2026}}
\label{fig:chuoikhoi}
\end{figure}}

% ---------- Luồng tương tác của ứng dụng phi tập trung ----------
\newcommand{\HinhLuongDApp}{%
\begin{figure}[htbp]
\centering
\begin{tikzpicture}[
    >=Stealth,
    hop/.style={draw, rounded corners=3pt, align=center, font=\footnotesize,
                text width=2.1cm, minimum height=1.1cm, fill=white},
    nhan/.style={font=\scriptsize, align=center}]
  \node[hop]                         (u)   at (0,0)     {Người dùng};
  \node[hop]                         (fe)  at (3.7,0)   {Giao diện web\\(Next.js,\\ethers.js)};
  \node[hop]                         (mm)  at (7.4,0)   {Ví\\MetaMask\\(giữ khóa riêng)};
  \node[hop, text width=2.6cm, minimum height=1.6cm, fill=gray!12]
                                     (net) at (11.3,0)  {Mạng Ethereum\\(Sepolia)\\nút RPC, EVM,\\hợp đồng thông minh};
  \node[hop]                         (be)  at (3.7,-3)  {Backend\\(Node.js, Express)};
  \node[hop]                         (db)  at (0,-3)    {CSDL off-chain\\(PostgreSQL)};

  \draw[->] (u) -- node[nhan, above] {thao tác} (fe);
  \draw[->] (fe) -- node[nhan, above] {yêu cầu ký\\(EIP-1193)} (mm);
  \draw[->] (mm) -- node[nhan, above] {giao dịch\\đã ký} (net);
  \draw[->, dashed] (fe.north) to[bend left=22]
      node[nhan, above] {đọc trạng thái on-chain (call)} (net.north);
  \draw[->] (net.south) |- node[nhan, pos=0.75, below] {sự kiện (Minted, Transfer, Retired\ldots)} (be.east);
  \draw[<->] (fe) -- node[nhan, right] {API REST} (be);
  \draw[<->] (be) -- node[nhan, above] {đọc/ghi} (db);
\end{tikzpicture}
\caption[Luồng tương tác giữa giao diện, ví MetaMask, mạng Ethereum và backend]{Luồng tương tác giữa giao diện, ví MetaMask, mạng Ethereum và backend \hinhtao[Claude]{20/10/2026}}
\label{fig:luongdapp}
\end{figure}}

\newcommand{\BangCongNghe}{%
\begin{table}[htbp]
\centering
\caption{Công nghệ và phiên bản sử dụng}
\label{tab:congnghe}
\small
\renewcommand{\arraystretch}{1.2}
\begin{tabularx}{\textwidth}{|p{3.2cm}|p{4.4cm}|X|}
\hline
\textbf{Thành phần} & \textbf{Công nghệ (phiên bản)} & \textbf{Vai trò}\\ \hline
Hợp đồng thông minh & Solidity 0.8.24, OpenZeppelin Contracts 5.1.0 & ERC-721, ERC-20, AccessControl, Pausable, ReentrancyGuard\\ \hline
Môi trường phát triển & Hardhat 2.29, hardhat-toolbox 5 (ethers 6, Mocha/Chai, solidity-coverage, gas reporter) & Biên dịch, kiểm thử, đo độ phủ và gas, triển khai\\ \hline
Mạng blockchain & Hardhat Network (cục bộ), Ethereum Sepolia & Kiểm thử và triển khai thử nghiệm\\ \hline
Backend & Node.js 22, Express 4.21, jsonwebtoken 9, multer & API REST, xác thực bằng chữ ký ví, tải hồ sơ\\ \hline
Cơ sở dữ liệu & PostgreSQL 16 (pg 8.13); pg-mem khi kiểm thử & Lưu dữ liệu off-chain theo mô hình CAD Trust\\ \hline
Đồng bộ sự kiện & ethers 6.13 (JsonRpcProvider, queryFilter) & Đọc sự kiện on-chain vào CSDL\\ \hline
Frontend & Next.js 14.2 (App Router), React 18.3, ethers 6.13 & Giao diện web, kết nối MetaMask\\ \hline
Ví & MetaMask & Ký đăng nhập và ký giao dịch\\ \hline
\end{tabularx}
\end{table}}

\newcommand{\HinhTrienKhai}{%
\begin{figure}[htbp]
\centering
\begin{adjustbox}{max width=\textwidth}
\begin{tikzpicture}[>=Stealth, font=\footnotesize,
    hop/.style={draw, rounded corners=3pt, align=center, fill=white, minimum height=1.1cm, text width=3.2cm},
    nhan/.style={font=\scriptsize, align=center}]
  \node[hop] (br)  at (0,0)    {\textbf{Trình duyệt}\\ giao diện Next.js\\ + MetaMask};
  \node[hop] (fe)  at (0,-3.2) {\textbf{Máy chủ frontend}\\ Next.js (Vercel)};
  \node[hop] (be)  at (5.2,0)  {\textbf{Máy chủ backend}\\ Express, port 4000\\ + bộ đồng bộ sự kiện};
  \node[hop] (db)  at (10.4,0) {\textbf{PostgreSQL}};
  \node[hop] (fs)  at (10.4,-3.2) {\textbf{Kho tệp}\\ uploads/ (hoặc IPFS)};
  \node[hop] (rpc) at (5.2,-3.2) {\textbf{Nút RPC}\\ (Infura/Alchemy)};
  \node[hop, fill=gray!10] (eth) at (5.2,-6.0) {\textbf{Ethereum Sepolia}\\ CarbonCreditNFT, CarbonPool};
  \draw[->] (fe) -- node[nhan, left] {tải trang\\ (HTTPS)} (br);
  \draw[<->] (br) -- node[nhan, above] {REST/JSON\\ + JWT} (be);
  \draw[<->] (be) -- node[nhan, above] {SQL} (db);
  \draw[->] (be) -- node[nhan, right] {lưu tệp} (fs);
  \draw[->] (be) -- node[nhan, right] {đọc sự kiện} (rpc);
  \draw[->] (br.south east) to[bend left=10] node[nhan, pos=0.35, right] {MetaMask ký\\ và gửi giao dịch} (rpc.north west);
  \draw[<->] (rpc) -- (eth);
\end{tikzpicture}
\end{adjustbox}
\caption[Sơ đồ triển khai hệ thống]{Sơ đồ triển khai hệ thống \hinhtao[Claude]{21/10/2026}}
\label{fig:trienkhai}
\end{figure}}

\newcommand{\HinhCauTrucMa}{%
\begin{figure}[htbp]
\centering
\footnotesize
\renewcommand{\arraystretch}{1.05}
\begin{tabular}{@{}l l@{}}
\texttt{DATN-CarbonCreditOnBlockChain/} & \\
\texttt{\quad contracts/} & \textit{Hardhat}\\
\texttt{\quad\quad contracts/CarbonCreditNFT.sol} & lô tín chỉ ERC-721: mint, retire\\
\texttt{\quad\quad contracts/CarbonPool.sol} & pool ERC-20: deposit, redeem, retire\\
\texttt{\quad\quad test/CarbonCredit.test.js} & 26 test case TC-01 đến TC-26\\
\texttt{\quad\quad scripts/deploy.js} & triển khai, cấp quyền\\
\texttt{\quad\quad scripts/demo.js} & kịch bản demo, đo gas\\
\texttt{\quad\quad scripts/e2e-indexer.js} & kiểm thử on-chain $\rightarrow$ CSDL\\
\texttt{\quad\quad scripts/smoke-api.js} & kiểm thử toàn hệ thống đang chạy\\
\texttt{\quad backend/} & \textit{Node.js + Express}\\
\texttt{\quad\quad db/schema.sql} & lược đồ CSDL (CAD Trust + ứng dụng)\\
\texttt{\quad\quad src/routes/*.js} & các nhóm endpoint API\\
\texttt{\quad\quad src/indexer.js} & đồng bộ sự kiện on-chain\\
\texttt{\quad\quad test/api.test.js} & 16 test case API-01 đến API-16\\
\texttt{\quad frontend/} & \textit{Next.js}\\
\texttt{\quad\quad app/} & các trang giao diện\\
\texttt{\quad\quad components/WalletProvider.js} & đăng nhập bằng ví\\
\texttt{\quad\quad lib/chain.js, lib/api.js} & gọi hợp đồng và API\\
\end{tabular}
\caption[Cấu trúc thư mục mã nguồn]{Cấu trúc thư mục mã nguồn \hinhtao[Claude]{21/10/2026}}
\label{fig:cautrucma}
\end{figure}}

% ============================================================
% BẢNG CHƯƠNG 5 — số liệu lấy trực tiếp từ kết quả chạy ngày 21/10/2026
%   (npx hardhat test / coverage / REPORT_GAS, npm test trong backend)
%   \BangTestHopDong : 26 test case hợp đồng               (tab:testhopdong)
%   \BangTestAPI     : 16 test case API                    (tab:testapi)
%   \BangDoPhu       : độ phủ mã hợp đồng                  (tab:dophu)
%   \BangGas         : chi phí gas các thao tác chính      (tab:gas)
%   \BangTinhToanKep : kịch bản chống tính toán kép        (tab:tinhtoankep)
% CHẠY LẠI và cập nhật số liệu nếu sửa hợp đồng. Cột "Yêu cầu" cần điền mã YC-xx
% sau khi bạn lập bảng yêu cầu chức năng ở mục \ref{subsec:ycchucnang}.
% ============================================================

\newcommand{\BangTestHopDong}{%
\begingroup\small
\renewcommand{\arraystretch}{1.12}
\begin{longtable}{|p{1.25cm}|p{8.8cm}|p{1.55cm}|p{1.55cm}|}
\caption{Kết quả kiểm thử đơn vị hợp đồng thông minh}\label{tab:testhopdong}\\
\hline
\textbf{Mã} & \textbf{Kịch bản kiểm thử} & \textbf{Yêu cầu} & \textbf{Kết quả}\\ \hline
\endfirsthead
\hline
\textbf{Mã} & \textbf{Kịch bản kiểm thử} & \textbf{Yêu cầu} & \textbf{Kết quả}\\ \hline
\endhead
\multicolumn{4}{|l|}{\textit{Phát hành (Minting)}}\\ \hline
TC-01 & VVB mint lô tín chỉ, lưu đúng thông tin và phát sự kiện CreditMinted & YC-?? & Đạt\\ \hline
TC-02 & Tài khoản không có VERIFIER\_ROLE không mint được & YC-?? & Đạt\\ \hline
TC-03 & Chặn dãy sê-ri chồng lấn trong cùng dự án và vintage (phát hành kép) & YC-?? & Đạt\\ \hline
TC-04 & Cho phép lô kế tiếp không chồng lấn, và cùng dãy sê-ri ở vintage khác & YC-?? & Đạt\\ \hline
TC-05 & Chặn dùng lại cùng một báo cáo thẩm định (phát hành kép) & YC-?? & Đạt\\ \hline
TC-06 & Từ chối số lượng 0 và sê-ri 0 & YC-?? & Đạt\\ \hline
\multicolumn{4}{|l|}{\textit{Chuyển nhượng và hủy (Retirement)}}\\ \hline
TC-07 & Chủ sở hữu chuyển nhượng NFT lô tín chỉ & YC-?? & Đạt\\ \hline
TC-08 & Hủy một phần lô, cấp chứng nhận với dãy sê-ri đúng & YC-?? & Đạt\\ \hline
TC-09 & Hủy hết lô thì NFT bị đốt, không chuyển nhượng hay hủy lại được (sử dụng kép) & YC-?? & Đạt\\ \hline
TC-10 & Không hủy quá số còn lại, không hủy 0, bắt buộc có người thụ hưởng & YC-?? & Đạt\\ \hline
TC-11 & Người không sở hữu không hủy được; người được approve thì hủy được & YC-?? & Đạt\\ \hline
TC-12 & retireFor chỉ dành cho RETIREMENT\_AGENT\_ROLE đang giữ NFT & YC-?? & Đạt\\ \hline
TC-13 & getBatch của lô không tồn tại bị từ chối & YC-?? & Đạt\\ \hline
TC-14 & Cơ quan quản lý tạm dừng thì chặn mint, chuyển nhượng và hủy & YC-?? & Đạt\\ \hline
TC-15 & Hỗ trợ các interface ERC-721 và AccessControl & YC-?? & Đạt\\ \hline
\multicolumn{4}{|l|}{\textit{Pool ERC-20 (giao dịch từng phần)}}\\ \hline
TC-16 & Nạp NFT vào pool nhận đúng số token, NFT bị khóa trong pool & YC-?? & Đạt\\ \hline
TC-17 & Nạp bằng safeTransferFrom trực tiếp; lô đã hủy một phần chỉ nhận phần còn lại & YC-?? & Đạt\\ \hline
TC-18 & Pool từ chối NFT không phải của hợp đồng tín chỉ & YC-?? & Đạt\\ \hline
TC-19 & Chuyển token ERC-20 từng phần cho người mua & YC-?? & Đạt\\ \hline
TC-20 & Người mua hủy từ pool, chứng nhận ghi người mua và dãy sê-ri; token đã đốt không dùng lại được (sử dụng kép) & YC-?? & Đạt\\ \hline
TC-21 & Bất biến: tổng cung token bằng tổng phần còn lại của các lô trong pool & YC-?? & Đạt\\ \hline
TC-22 & Đổi lại (redeem) lô cần đốt đủ token bằng phần còn lại & YC-?? & Đạt\\ \hline
TC-23 & Thao tác với lô không nằm trong pool bị từ chối & YC-?? & Đạt\\ \hline
TC-24 & NFT trong pool không thể bị chủ cũ bán riêng (bán kép) & YC-?? & Đạt\\ \hline
TC-25 & Hủy từ pool quá phần còn lại của lô bị từ chối, số dư không đổi & YC-?? & Đạt\\ \hline
TC-26 & Nạp lô đã hủy hết (không còn tồn tại) bị từ chối & YC-?? & Đạt\\ \hline
\end{longtable}
\endgroup}

\newcommand{\BangTestAPI}{%
\begingroup\small
\renewcommand{\arraystretch}{1.12}
\begin{longtable}{|p{1.35cm}|p{8.7cm}|p{1.55cm}|p{1.55cm}|}
\caption{Kết quả kiểm thử API và bộ đồng bộ sự kiện}\label{tab:testapi}\\
\hline
\textbf{Mã} & \textbf{Kịch bản kiểm thử} & \textbf{Yêu cầu} & \textbf{Kết quả}\\ \hline
\endfirsthead
\hline
\textbf{Mã} & \textbf{Kịch bản kiểm thử} & \textbf{Yêu cầu} & \textbf{Kết quả}\\ \hline
\endhead
API-01 & Đăng nhập bằng chữ ký ví, tài khoản mới tự chọn vai trò & YC-?? & Đạt\\ \hline
API-02 & Chữ ký sai hoặc dùng lại nonce bị từ chối & YC-?? & Đạt\\ \hline
API-03 & Không tự đăng ký được vai trò VERIFIER; ADMIN cấp được & YC-?? & Đạt\\ \hline
API-04 & Endpoint quản trị từ chối người không phải ADMIN và người chưa đăng nhập & YC-?? & Đạt\\ \hline
API-05 & Chủ dự án đăng ký dự án; trùng mã dự án bị từ chối & YC-?? & Đạt\\ \hline
API-06 & Gửi yêu cầu cấp tín chỉ cần báo cáo giám sát đã tải lên & YC-?? & Đạt\\ \hline
API-07 & Chỉ VVB duyệt được; duyệt xong trả về tham số mint với sê-ri đề xuất & YC-?? & Đạt\\ \hline
API-08 & Báo cáo thẩm định đã dùng không duyệt được yêu cầu khác (phát hành kép, off-chain) & YC-?? & Đạt\\ \hline
API-09 & Sự kiện CreditMinted tạo đợt cấp, khối tín chỉ và chuyển yêu cầu sang MINTED & YC-?? & Đạt\\ \hline
API-10 & Sự kiện trùng (cùng tx\_hash, log\_index) không bị ghi hai lần & YC-?? & Đạt\\ \hline
API-11 & Sự kiện CreditRetired tách khối và tạo chứng nhận tra cứu được & YC-?? & Đạt\\ \hline
API-12 & Tra cứu theo số sê-ri cho biết sê-ri đã hủy hay còn hiệu lực & YC-?? & Đạt\\ \hline
API-13 & Dashboard so sánh tín chỉ đã hủy với phần phát thải vượt hạn ngạch & YC-?? & Đạt\\ \hline
API-14 & Hủy hết phần còn lại thì không còn khối Held; thống kê khớp & YC-?? & Đạt\\ \hline
API-15 & Xuất dữ liệu dự án theo cấu trúc CAD Trust & YC-?? & Đạt\\ \hline
API-16 & Metadata NFT và tài liệu công khai cho phép đối chiếu mã băm & YC-?? & Đạt\\ \hline
E2E-01 & Bộ đồng bộ đọc sự kiện thật trên mạng Hardhat (mint, nạp pool, hủy); chạy lại không ghi trùng; phần còn lại on-chain khớp CSDL & YC-?? & Đạt\\ \hline
E2E-02 & Kiểm thử toàn hệ thống đang chạy: đăng nhập ví, đăng ký dự án, yêu cầu, VVB duyệt và mint, nạp pool, bán, hủy, tra cứu, dashboard & YC-?? & Đạt\\ \hline
\end{longtable}
\endgroup}

\newcommand{\BangDoPhu}{%
\begin{table}[htbp]
\centering
\caption{Độ phủ mã của bộ kiểm thử hợp đồng (solidity-coverage)}
\label{tab:dophu}
\begin{tabular}{|l|r|r|r|r|}
\hline
\textbf{Tệp} & \textbf{Câu lệnh} & \textbf{Nhánh} & \textbf{Hàm} & \textbf{Dòng}\\ \hline
CarbonCreditNFT.sol & 100\% & 94,44\% & 100\% & 100\%\\ \hline
CarbonPool.sol & 100\% & 75,00\% & 100\% & 100\%\\ \hline
\textbf{Tổng} & \textbf{100\%} & \textbf{88,46\%} & \textbf{100\%} & \textbf{100\%}\\ \hline
\end{tabular}
\end{table}}

\newcommand{\BangGas}{%
\begin{table}[htbp]
\centering
\caption{Chi phí gas của các thao tác chính (Solidity 0.8.24, optimizer 200 runs)}
\label{tab:gas}
\small
\begin{tabular}{|l|l|r|r|r|}
\hline
\textbf{Hợp đồng} & \textbf{Thao tác} & \textbf{Nhỏ nhất} & \textbf{Lớn nhất} & \textbf{Trung bình}\\ \hline
CarbonCreditNFT & Triển khai & -- & -- & 2.431.505\\ \hline
CarbonPool & Triển khai & -- & -- & 1.167.804\\ \hline
CarbonCreditNFT & mintBatch & 260.483 & 294.683 & 288.105\\ \hline
CarbonCreditNFT & retire (trực tiếp) & 206.215 & 261.001 & 243.049\\ \hline
CarbonCreditNFT & transferFrom & -- & -- & 57.310\\ \hline
CarbonCreditNFT & safeTransferFrom vào pool & -- & -- & 182.071\\ \hline
CarbonPool & deposit & 116.482 & 184.882 & 157.522\\ \hline
CarbonPool & transfer (PCC) & 29.683 & 51.595 & 47.208\\ \hline
CarbonPool & retire (từ pool) & 262.346 & 282.681 & 274.041\\ \hline
CarbonPool & redeem & -- & -- & 84.162\\ \hline
\end{tabular}
\end{table}}

\newcommand{\BangTinhToanKep}{%
\begin{table}[htbp]
\centering
\caption{Kịch bản chống tính toán kép và cơ chế chặn}
\label{tab:tinhtoankep}
\small
\renewcommand{\arraystretch}{1.2}
\begin{tabularx}{\textwidth}{|p{2.6cm}|X|X|p{1.9cm}|}
\hline
\textbf{Dạng} & \textbf{Kịch bản tấn công} & \textbf{Cơ chế chặn} & \textbf{Test}\\ \hline
Phát hành kép & VVB mint lô có dãy sê-ri chồng lên lô đã cấp của cùng dự án và vintage & nextAvailableSerial: sê-ri phải tăng dần, lỗi SerialOverlap & TC-03\\ \hline
Phát hành kép & Dùng lại một báo cáo thẩm định để cấp lần hai (kể cả cho dự án khác) & documentUsed on-chain (DocumentAlreadyUsed) và verifications.document\_hash unique off-chain & TC-05, API-08\\ \hline
Sử dụng kép & Hủy lại hoặc chuyển nhượng tín chỉ đã hủy & Trừ phần đã hủy khỏi lô; hủy hết thì đốt NFT & TC-09, TC-10\\ \hline
Sử dụng kép & Dùng lại token PCC đã đốt khi hủy từ pool & \_burn trước khi gọi retireFor; số dư không đủ thì giao dịch bị hoàn tác & TC-20\\ \hline
Bán kép & Chủ cũ bán riêng NFT đã nạp vào pool (đã nhận token tương ứng) & NFT thuộc sở hữu của pool, chỉ ra qua redeem khi đốt đủ token & TC-24, TC-21\\ \hline
Tuyên bố kép (trong hệ thống) & Hai bên cùng nhận đã bù trừ bằng một dãy sê-ri & Mỗi chứng nhận ghi dãy sê-ri và người thụ hưởng duy nhất; tra cứu công khai theo sê-ri & TC-08, API-12\\ \hline
\end{tabularx}
\end{table}}


% Tổng quan (Normal, không tiêu đề): chương trình bày các công nghệ đồ án dùng để hiện thực thiết kế ở
% Chương 3, kết quả cài đặt, kiểm thử và triển khai trên Sepolia.
% [VIẾT SAU] \canviet{Đoạn tổng quan chương}

% ============ PHẦN A — CÔNG NGHỆ SỬ DỤNG (từ Chương 2 cũ) ============
% [VIẾT SAU] Toàn bộ các mục dưới đây chưa viết nội dung — bỏ dấu % để khôi phục.
% % ------------------------------------------------------------
% \section{Blockchain và nền tảng Ethereum}
% \label{sec:blockchain}
% \label{sec:ethereum}
%
% \subsection*{Blockchain}
% \label{subsec:cautrucblock}
% \label{subsec:dongthuan}
% % Ý cần có (ngắn gọn, chỉ phần đồ án dùng): khối, mã băm khối trước, sổ cái chỉ ghi thêm (ledger) ->
% % vì sao dữ liệu on-chain khó sửa, phục vụ yêu cầu toàn vẹn ở mục \ref{sec:phichucnang} và VĐ-2;
% % đồng thuận PoW/PoS, Ethereum chuyển sang PoS (The Merge 15/9/2022).
% % Trích: \cite{nakamoto2008bitcoin,zheng2017overview,ethereum2024pos}. Nhắc \autoref{fig:chuoikhoi} trước hình.
% % [VIẾT SAU] \canviet{Blockchain}
% %
% % \HinhChuoiKhoi
%
% \subsection*{Lựa chọn nền tảng blockchain}
% \label{subsec:phanloai}
% % Bảng so sánh Ethereum (public) với Hyperledger Fabric (permissioned, mà Vaccargiu dùng) và Chia
% % (CAD Trust dùng): mô hình truy cập, đồng thuận, ngôn ngữ smart contract, chuẩn token, ví, testnet, chi phí.
% % Lý do chọn Ethereum: phiếu giao nhiệm vụ; chuẩn ERC-721/ERC-20; MetaMask; testnet công khai để bên thứ ba
% % tự kiểm tra (VĐ-2, NC-07). Nhược điểm: phí gas, dữ liệu công khai -> xử lý ở yêu cầu quyền riêng tư.
% % [VIẾT SAU] \canviet{Lựa chọn nền tảng blockchain}
%
% \subsection*{Tài khoản, giao dịch và phí gas}
% \label{subsec:taikhoan}
% \label{subsec:evm}
% % Ý cần có: EOA và contract account; khóa riêng, địa chỉ ví; giao dịch ký số, nonce; EVM, gas, EIP-1559.
% % Đây là cơ sở cho phần đo gas ở Chương 3. Trích: \cite{buterin2014ethereum,wood2014yellow,eip1559}.
% % [VIẾT SAU] \canviet{Tài khoản, giao dịch và phí gas}
%
% \subsection*{Testnet Sepolia và Etherscan}
% \label{subsec:sepolia}
% % Ý cần có: mainnet và testnet; Sepolia, faucet; Etherscan để tra giao dịch, xác minh (verify) mã hợp đồng.
% % Lý do: phạm vi nguyên mẫu (mục \ref{sec:muctieu}), không tốn phí thật.
% % [VIẾT SAU] \canviet{Testnet Sepolia}
%
% % ------------------------------------------------------------
% \section{Smart contract và công cụ phát triển}
% \label{sec:smartcontract}
%
% \subsection*{Smart contract và ngôn ngữ Solidity}
% \label{subsec:khainiemsc}
% \label{subsec:solidity}
% % Ý cần có: định nghĩa (Szabo); tự thực thi, không sửa được mã sau triển khai; event và log (dùng cho
% % bộ đồng bộ sự kiện); cấu trúc contract (state variable, function, modifier, event), kế thừa.
% % Yêu cầu phục vụ: NC-03, NC-06 (ràng buộc sê-ri, không hủy hai lần thành ràng buộc kỹ thuật).
% % Trích: \cite{szabo1997formalizing,buterin2014ethereum,solidity2024docs}.
% % [VIẾT SAU] \canviet{Smart contract và Solidity}
%
% \subsection*{Thư viện OpenZeppelin và Hardhat}
% \label{subsec:openzeppelin}
% % OpenZeppelin 5: ERC721, ERC20, AccessControl (VERIFIER_ROLE...), Pausable, ReentrancyGuard -> NC-09 và
% % yêu cầu bảo mật. Hardhat: biên dịch, Hardhat Network, test Mocha/Chai, coverage, gas reporter.
% % Phương án thay thế: Foundry, Truffle (đã ngừng phát triển); lý do chọn Hardhat (JavaScript cùng ngôn ngữ backend).
% % Trích: \cite{openzeppelin2024contracts}.
% % [VIẾT SAU] \canviet{OpenZeppelin và Hardhat}
%
% \subsection*{Các lỗ hổng thường gặp của smart contract}
% \label{subsec:baomat}
% % Reentrancy, kiểm soát truy cập sai, tràn số (Solidity >= 0.8 tự kiểm tra), tx.origin. Mỗi lỗ hổng
% % 1–2 câu và cách phòng trong đồ án; công cụ quét Slither. Trích: \cite{atzei2017survey}.
% % [VIẾT SAU] \canviet{Lỗ hổng smart contract}
%
% % ------------------------------------------------------------
% \section{Các chuẩn token}
% \label{sec:token-standards}
%
% \subsection*{ERC-20 và ERC-721}
% \label{subsec:erc20}
% \label{subsec:erc721}
% % ERC-20: fungible token, totalSupply, balanceOf, transfer, approve, transferFrom, allowance \cite{eip20}.
% % ERC-721: NFT, ownerOf, safeTransferFrom, tokenURI, onERC721Received \cite{eip721}.
% % [VIẾT SAU] \canviet{ERC-20 và ERC-721}
%
% \subsection*{Lựa chọn chuẩn token cho tín chỉ các-bon}
% \label{subsec:chontoken}
% % Bảng so sánh ERC-20 / ERC-721 / ERC-1155 theo: giữ metadata riêng từng lô, chia nhỏ, thanh khoản,
% % hỗ trợ ví. Kết luận: mỗi lô là một NFT ERC-721 (giữ dự án, vintage, dãy sê-ri, NC-03), pool ERC-20
% % cho giao dịch từng phần (NC-05); kế thừa mô hình Batch NFT/TCO2/BCT của Toucan (mục \ref{sec:khaosatsobo}).
% % [VIẾT SAU] \canviet{Lựa chọn chuẩn token}
%
% % ------------------------------------------------------------
% \section{Ứng dụng web giao tiếp với blockchain}
% \label{sec:dapp}
%
% \subsection*{Kiến trúc dApp}
% \label{subsec:kientrucdapp}
% % Luồng một giao dịch từ giao diện tới smart contract (giải thích \autoref{fig:luongdapp}).
% % [VIẾT SAU] \canviet{Kiến trúc dApp}
% %
% % \HinhLuongDApp
%
% \subsection*{Ví MetaMask, ethers.js và Next.js}
% \label{subsec:metamask}
% \label{subsec:ethers}
% % MetaMask giữ khóa riêng, EIP-1193 (window.ethereum) \cite{eip1193}; ethers.js: Provider, Signer,
% % Contract, ABI, call và send \cite{ethers2024docs}; Next.js cho giao diện. Yêu cầu phục vụ: tính dễ dùng,
% % người dùng tự ký giao dịch (backend không giữ khóa). Phương án thay thế: web3.js, viem; React thuần.
% % [VIẾT SAU] \canviet{MetaMask, ethers.js và Next.js}
%
% % ------------------------------------------------------------
% \section{Công nghệ phía off-chain}
% \label{sec:offchain-congnghe}
% \label{sec:onoffchain-lythuyet}
% % Node.js + Express (API RESTful), PostgreSQL (dữ liệu theo CAD Trust V2.0, NC-10), JWT, xác thực bằng
% % chữ ký ví. Tiêu chí chọn dữ liệu on-chain/off-chain theo \cite{xu2017taxonomy}; neo mã băm keccak256 của hồ sơ;
% % IPFS \cite{benet2014ipfs} là hướng mở rộng (hiện lưu tệp cục bộ). Phương án thay thế: MongoDB; lý do chọn
% % CSDL quan hệ (mô hình CAD Trust có quan hệ rõ, giống UNDP dùng PostgreSQL).
% % [VIẾT SAU] \canviet{Công nghệ off-chain}


% ============ PHẦN B — XÂY DỰNG, KIỂM THỬ, TRIỂN KHAI (từ Chương 3 cũ) ============
% ------------------------------------------------------------
\section{Xây dựng ứng dụng}
\label{sec:xaydung}

\subsection*{Thư viện và công cụ sử dụng}
\label{sec:moitruong}
% Mẫu SoICT: bảng công cụ có phiên bản (và URL). Giải thích \autoref{tab:congnghe}; lý do chọn công cụ đã nêu
% ở Chương 2 nên chỉ nhắc lại ngắn. Cấu trúc mã nguồn theo \autoref{fig:cautrucma}.
% [VIẾT SAU] \canviet{Thư viện và công cụ sử dụng}

\BangCongNghe

\HinhCauTrucMa

\subsection*{Kết quả đạt được}
\label{subsec:ketquadatduoc}
\label{sec:caidat-hopdong}
\label{sec:caidat-backend}
% Mẫu SoICT: mô tả các sản phẩm đóng gói, thành phần, vai trò; bảng thống kê số dòng mã, số tệp, số hợp đồng,
% số endpoint, dung lượng.
% Nội dung cài đặt từ Chương 4 cũ (chỉ tóm lược, chi tiết đóng góp ở Chương 4):
%  - CarbonCreditNFT: struct CreditBatch (projectId, vintage, serialFrom, quantity, retired, documentHash,
%    verifier, issuedAt); mintBatch: seriesKey = keccak256(projectId, vintage), serialFrom >= nextAvailableSerial,
%    documentUsed[documentHash]; retire hủy từ đầu dãy sê-ri còn lại; retireFor cho pool (RETIREMENT_AGENT_ROLE);
%    Pausable ghi đè _update.
%  - CarbonPool: decimals = 0; nạp qua onERC721Received; redeem đốt remainingOf(tokenId); retire đốt token rồi
%    gọi retireFor; ReentrancyGuard.
%  - Phân quyền: DEFAULT_ADMIN_ROLE và vai trò phát hành ISSUER_ROLE (cơ quan quản lý, mỗi tỉnh một tài khoản),
%    VERIFIER_ROLE (VVB, chỉ phê duyệt off-chain), RETIREMENT_AGENT_ROLE (pool). [CẦN SỬA MÃ] thêm ISSUER_ROLE.
%  - Backend: db/schema.sql từ ERD; xác thực nonce -> personal_sign -> ethers.verifyMessage -> JWT 8 giờ;
%    API theo \autoref{tab:api} (backend không giữ khóa riêng); indexer quét queryFilter, khóa (tx_hash, log_index),
%    chưa xử lý reorg.
% [VIẾT SAU] \canviet{Kết quả đạt được}

% \BangAPI   % [17/11/2026] đã dùng ở Chương 3, mục Thiết kế API

\subsection*{Minh họa các chức năng chính}
\label{subsec:minhhoa}
\label{sec:caidat-frontend}
% Mỗi màn hình 1 ảnh chụp + giải thích ngắn: danh sách dự án, chi tiết dự án, tra cứu sê-ri, chứng nhận hủy,
% dự án của tôi, thẩm định (VVB), ví tín chỉ, dashboard, quản trị. Chụp khi chạy thật (README "Chạy thử cục bộ").
% [VIẾT SAU] \canviet{Minh họa các chức năng chính}

% ------------------------------------------------------------
\section{Kiểm thử}
\label{sec:kiemthu}
\label{sec:pp-kiemthu}
\label{sec:kq-chucnang}
% Mẫu SoICT: 2–3 trang; thiết kế ca kiểm thử cho 2–3 chức năng quan trọng nhất (phát hành, hủy từ pool, tra cứu),
% nêu kỹ thuật kiểm thử; tổng kết số ca và kết quả; giải thích ca không đạt. Ca kiểm thử khác đưa vào Phụ lục.
% Ba mức: kiểm thử đơn vị hợp đồng trên Hardhat Network (Mocha/Chai, loadFixture); kiểm thử API với pg-mem và
% supertest; kiểm thử tích hợp e2e-indexer và smoke-api. Ghi cấu hình máy, phiên bản Node, Hardhat.
% Nhận xét độ phủ: nhánh của CarbonPool 75% vì ... (xem coverage/index.html).
% Chi phí gas: quy đổi ETH/USD theo giá gas tại thời điểm đo (ghi nguồn, ngày); vì sao mintBatch, retire tốn gas.
% Kịch bản chống tính toán kép (\BangTinhToanKep) đặt ở Chương 4, mục \ref{sec:gp-tinhtoankep}.
% [VIẾT SAU] \canviet{Kiểm thử}

\BangTestHopDong

\BangTestAPI

\BangDoPhu

\BangGas

% ------------------------------------------------------------
\section{Triển khai}
\label{sec:trienkhai}
\label{sec:kq-sepolia}
% Giải thích \autoref{fig:trienkhai}: Sepolia + RPC, backend + PostgreSQL, frontend; địa chỉ hợp đồng,
% link Etherscan (đã verify), vài giao dịch mẫu (mint, retire), ảnh chụp.
% [VIẾT SAU] \canviet{Triển khai}

\HinhTrienKhai


% ------------------------------------------------------------
% Kết chương (Normal, không tiêu đề): tóm tắt bộ công nghệ (bảng Công nghệ | Phiên bản | Vai trò | Yêu cầu
% phục vụ), kết quả cài đặt, kiểm thử, triển khai; dẫn sang Chương 5 (giải pháp và đóng góp).
% [VIẾT SAU] \canviet{Kết chương 4}
 để khôi phục.
% Bảng API (\BangAPI) đã dùng ở Chương 3 (Thiết kế API), không đặt lại ở đây.
%
% Hướng dẫn của mẫu SoICT cho phần công nghệ: không quá 10 trang; chỉ giới thiệu công nghệ đồ án thực sự
% dùng; với mỗi công nghệ nêu (i) yêu cầu nó phục vụ (NC-xx, VĐ-x, yêu cầu phi chức năng ở Chương 2),
% (ii) phương án thay thế, (iii) lý do chọn.
% Văn phong: xem ghi chú đầu Chap1.tex. [VIẾT SAU] toàn chương.
% ============================================================

\chapter{CÔNG NGHỆ, XÂY DỰNG VÀ KIỂM THỬ HỆ THỐNG}
\chaptermark{CÔNG NGHỆ, XÂY DỰNG VÀ KIỂM THỬ}
\label{chap:congnghe}
\label{chap:cosolythuyet}
\label{chap:xaydung}
\label{chap:ketqua}
% ============================================================
% HÌNH VÀ BẢNG CHƯƠNG 4 (CÔNG NGHỆ, XÂY DỰNG VÀ KIỂM THỬ) — gọi bằng macro trong Chap4.tex
%   \HinhChuoiKhoi  : cấu trúc liên kết các khối           (fig:chuoikhoi)   [chuyển từ figures_chap2 ngày 17/11/2026]
%   \HinhLuongDApp  : luồng tương tác của dApp              (fig:luongdapp)   [chuyển từ figures_chap2 ngày 17/11/2026]
%   \BangCongNghe   : công nghệ và phiên bản sử dụng      (tab:congnghe)
%   \HinhTrienKhai  : sơ đồ triển khai hệ thống            (fig:trienkhai)
%   \HinhCauTrucMa  : cấu trúc thư mục mã nguồn            (fig:cautrucma)
% Phiên bản lấy từ package.json / hardhat.config.js tại thời điểm 21/10/2026.
% ============================================================

% ---------- Cấu trúc liên kết các khối ----------
\newcommand{\HinhChuoiKhoi}{%
\begin{figure}[htbp]
\centering
\begin{tikzpicture}[
    >=Stealth,
    truong/.style={draw, minimum width=3.3cm, minimum height=0.55cm,
                   font=\footnotesize, fill=white, inner sep=1pt},
    tieude/.style={truong, fill=gray!20, font=\footnotesize\bfseries}]
  \foreach \i/\x/\ten in {0/0/N-1, 1/4.6/N, 2/9.2/N+1} {
    \node[tieude] (h\i) at (\x,0) {Khối \ten};
    \node[truong, below=0pt of h\i] (p\i) {Mã băm khối trước};
    \node[truong, below=0pt of p\i] (t\i) {Dấu thời gian};
    \node[truong, below=0pt of t\i] (m\i) {Gốc Merkle};
    \node[truong, below=0pt of m\i, minimum height=1cm, align=center] (g\i)
        {Danh sách\\giao dịch};
    \node[draw, thick, rounded corners=2pt, inner sep=3pt, fit=(h\i)(g\i)] (b\i) {};
  }
  \draw[->, thick] (p1.west) -- (b0.east |- p1.west);
  \draw[->, thick] (p2.west) -- (b1.east |- p2.west);
  \node[font=\large] at (-2.35,-1.4) {$\cdots$};
  \node[font=\large] at (11.55,-1.4) {$\cdots$};
\end{tikzpicture}
\caption[Cấu trúc liên kết giữa các khối trong chuỗi khối]{Cấu trúc liên kết giữa các khối trong chuỗi khối \hinhtao[Claude]{20/10/2026}}
\label{fig:chuoikhoi}
\end{figure}}

% ---------- Luồng tương tác của ứng dụng phi tập trung ----------
\newcommand{\HinhLuongDApp}{%
\begin{figure}[htbp]
\centering
\begin{tikzpicture}[
    >=Stealth,
    hop/.style={draw, rounded corners=3pt, align=center, font=\footnotesize,
                text width=2.1cm, minimum height=1.1cm, fill=white},
    nhan/.style={font=\scriptsize, align=center}]
  \node[hop]                         (u)   at (0,0)     {Người dùng};
  \node[hop]                         (fe)  at (3.7,0)   {Giao diện web\\(Next.js,\\ethers.js)};
  \node[hop]                         (mm)  at (7.4,0)   {Ví\\MetaMask\\(giữ khóa riêng)};
  \node[hop, text width=2.6cm, minimum height=1.6cm, fill=gray!12]
                                     (net) at (11.3,0)  {Mạng Ethereum\\(Sepolia)\\nút RPC, EVM,\\hợp đồng thông minh};
  \node[hop]                         (be)  at (3.7,-3)  {Backend\\(Node.js, Express)};
  \node[hop]                         (db)  at (0,-3)    {CSDL off-chain\\(PostgreSQL)};

  \draw[->] (u) -- node[nhan, above] {thao tác} (fe);
  \draw[->] (fe) -- node[nhan, above] {yêu cầu ký\\(EIP-1193)} (mm);
  \draw[->] (mm) -- node[nhan, above] {giao dịch\\đã ký} (net);
  \draw[->, dashed] (fe.north) to[bend left=22]
      node[nhan, above] {đọc trạng thái on-chain (call)} (net.north);
  \draw[->] (net.south) |- node[nhan, pos=0.75, below] {sự kiện (Minted, Transfer, Retired\ldots)} (be.east);
  \draw[<->] (fe) -- node[nhan, right] {API REST} (be);
  \draw[<->] (be) -- node[nhan, above] {đọc/ghi} (db);
\end{tikzpicture}
\caption[Luồng tương tác giữa giao diện, ví MetaMask, mạng Ethereum và backend]{Luồng tương tác giữa giao diện, ví MetaMask, mạng Ethereum và backend \hinhtao[Claude]{20/10/2026}}
\label{fig:luongdapp}
\end{figure}}

\newcommand{\BangCongNghe}{%
\begin{table}[htbp]
\centering
\caption{Công nghệ và phiên bản sử dụng}
\label{tab:congnghe}
\small
\renewcommand{\arraystretch}{1.2}
\begin{tabularx}{\textwidth}{|p{3.2cm}|p{4.4cm}|X|}
\hline
\textbf{Thành phần} & \textbf{Công nghệ (phiên bản)} & \textbf{Vai trò}\\ \hline
Hợp đồng thông minh & Solidity 0.8.24, OpenZeppelin Contracts 5.1.0 & ERC-721, ERC-20, AccessControl, Pausable, ReentrancyGuard\\ \hline
Môi trường phát triển & Hardhat 2.29, hardhat-toolbox 5 (ethers 6, Mocha/Chai, solidity-coverage, gas reporter) & Biên dịch, kiểm thử, đo độ phủ và gas, triển khai\\ \hline
Mạng blockchain & Hardhat Network (cục bộ), Ethereum Sepolia & Kiểm thử và triển khai thử nghiệm\\ \hline
Backend & Node.js 22, Express 4.21, jsonwebtoken 9, multer & API REST, xác thực bằng chữ ký ví, tải hồ sơ\\ \hline
Cơ sở dữ liệu & PostgreSQL 16 (pg 8.13); pg-mem khi kiểm thử & Lưu dữ liệu off-chain theo mô hình CAD Trust\\ \hline
Đồng bộ sự kiện & ethers 6.13 (JsonRpcProvider, queryFilter) & Đọc sự kiện on-chain vào CSDL\\ \hline
Frontend & Next.js 14.2 (App Router), React 18.3, ethers 6.13 & Giao diện web, kết nối MetaMask\\ \hline
Ví & MetaMask & Ký đăng nhập và ký giao dịch\\ \hline
\end{tabularx}
\end{table}}

\newcommand{\HinhTrienKhai}{%
\begin{figure}[htbp]
\centering
\begin{adjustbox}{max width=\textwidth}
\begin{tikzpicture}[>=Stealth, font=\footnotesize,
    hop/.style={draw, rounded corners=3pt, align=center, fill=white, minimum height=1.1cm, text width=3.2cm},
    nhan/.style={font=\scriptsize, align=center}]
  \node[hop] (br)  at (0,0)    {\textbf{Trình duyệt}\\ giao diện Next.js\\ + MetaMask};
  \node[hop] (fe)  at (0,-3.2) {\textbf{Máy chủ frontend}\\ Next.js (Vercel)};
  \node[hop] (be)  at (5.2,0)  {\textbf{Máy chủ backend}\\ Express, port 4000\\ + bộ đồng bộ sự kiện};
  \node[hop] (db)  at (10.4,0) {\textbf{PostgreSQL}};
  \node[hop] (fs)  at (10.4,-3.2) {\textbf{Kho tệp}\\ uploads/ (hoặc IPFS)};
  \node[hop] (rpc) at (5.2,-3.2) {\textbf{Nút RPC}\\ (Infura/Alchemy)};
  \node[hop, fill=gray!10] (eth) at (5.2,-6.0) {\textbf{Ethereum Sepolia}\\ CarbonCreditNFT, CarbonPool};
  \draw[->] (fe) -- node[nhan, left] {tải trang\\ (HTTPS)} (br);
  \draw[<->] (br) -- node[nhan, above] {REST/JSON\\ + JWT} (be);
  \draw[<->] (be) -- node[nhan, above] {SQL} (db);
  \draw[->] (be) -- node[nhan, right] {lưu tệp} (fs);
  \draw[->] (be) -- node[nhan, right] {đọc sự kiện} (rpc);
  \draw[->] (br.south east) to[bend left=10] node[nhan, pos=0.35, right] {MetaMask ký\\ và gửi giao dịch} (rpc.north west);
  \draw[<->] (rpc) -- (eth);
\end{tikzpicture}
\end{adjustbox}
\caption[Sơ đồ triển khai hệ thống]{Sơ đồ triển khai hệ thống \hinhtao[Claude]{21/10/2026}}
\label{fig:trienkhai}
\end{figure}}

\newcommand{\HinhCauTrucMa}{%
\begin{figure}[htbp]
\centering
\footnotesize
\renewcommand{\arraystretch}{1.05}
\begin{tabular}{@{}l l@{}}
\texttt{DATN-CarbonCreditOnBlockChain/} & \\
\texttt{\quad contracts/} & \textit{Hardhat}\\
\texttt{\quad\quad contracts/CarbonCreditNFT.sol} & lô tín chỉ ERC-721: mint, retire\\
\texttt{\quad\quad contracts/CarbonPool.sol} & pool ERC-20: deposit, redeem, retire\\
\texttt{\quad\quad test/CarbonCredit.test.js} & 26 test case TC-01 đến TC-26\\
\texttt{\quad\quad scripts/deploy.js} & triển khai, cấp quyền\\
\texttt{\quad\quad scripts/demo.js} & kịch bản demo, đo gas\\
\texttt{\quad\quad scripts/e2e-indexer.js} & kiểm thử on-chain $\rightarrow$ CSDL\\
\texttt{\quad\quad scripts/smoke-api.js} & kiểm thử toàn hệ thống đang chạy\\
\texttt{\quad backend/} & \textit{Node.js + Express}\\
\texttt{\quad\quad db/schema.sql} & lược đồ CSDL (CAD Trust + ứng dụng)\\
\texttt{\quad\quad src/routes/*.js} & các nhóm endpoint API\\
\texttt{\quad\quad src/indexer.js} & đồng bộ sự kiện on-chain\\
\texttt{\quad\quad test/api.test.js} & 16 test case API-01 đến API-16\\
\texttt{\quad frontend/} & \textit{Next.js}\\
\texttt{\quad\quad app/} & các trang giao diện\\
\texttt{\quad\quad components/WalletProvider.js} & đăng nhập bằng ví\\
\texttt{\quad\quad lib/chain.js, lib/api.js} & gọi hợp đồng và API\\
\end{tabular}
\caption[Cấu trúc thư mục mã nguồn]{Cấu trúc thư mục mã nguồn \hinhtao[Claude]{21/10/2026}}
\label{fig:cautrucma}
\end{figure}}

% ============================================================
% BẢNG CHƯƠNG 5 — số liệu lấy trực tiếp từ kết quả chạy ngày 21/10/2026
%   (npx hardhat test / coverage / REPORT_GAS, npm test trong backend)
%   \BangTestHopDong : 26 test case hợp đồng               (tab:testhopdong)
%   \BangTestAPI     : 16 test case API                    (tab:testapi)
%   \BangDoPhu       : độ phủ mã hợp đồng                  (tab:dophu)
%   \BangGas         : chi phí gas các thao tác chính      (tab:gas)
%   \BangTinhToanKep : kịch bản chống tính toán kép        (tab:tinhtoankep)
% CHẠY LẠI và cập nhật số liệu nếu sửa hợp đồng. Cột "Yêu cầu" cần điền mã YC-xx
% sau khi bạn lập bảng yêu cầu chức năng ở mục \ref{subsec:ycchucnang}.
% ============================================================

\newcommand{\BangTestHopDong}{%
\begingroup\small
\renewcommand{\arraystretch}{1.12}
\begin{longtable}{|p{1.25cm}|p{8.8cm}|p{1.55cm}|p{1.55cm}|}
\caption{Kết quả kiểm thử đơn vị hợp đồng thông minh}\label{tab:testhopdong}\\
\hline
\textbf{Mã} & \textbf{Kịch bản kiểm thử} & \textbf{Yêu cầu} & \textbf{Kết quả}\\ \hline
\endfirsthead
\hline
\textbf{Mã} & \textbf{Kịch bản kiểm thử} & \textbf{Yêu cầu} & \textbf{Kết quả}\\ \hline
\endhead
\multicolumn{4}{|l|}{\textit{Phát hành (Minting)}}\\ \hline
TC-01 & VVB mint lô tín chỉ, lưu đúng thông tin và phát sự kiện CreditMinted & YC-?? & Đạt\\ \hline
TC-02 & Tài khoản không có VERIFIER\_ROLE không mint được & YC-?? & Đạt\\ \hline
TC-03 & Chặn dãy sê-ri chồng lấn trong cùng dự án và vintage (phát hành kép) & YC-?? & Đạt\\ \hline
TC-04 & Cho phép lô kế tiếp không chồng lấn, và cùng dãy sê-ri ở vintage khác & YC-?? & Đạt\\ \hline
TC-05 & Chặn dùng lại cùng một báo cáo thẩm định (phát hành kép) & YC-?? & Đạt\\ \hline
TC-06 & Từ chối số lượng 0 và sê-ri 0 & YC-?? & Đạt\\ \hline
\multicolumn{4}{|l|}{\textit{Chuyển nhượng và hủy (Retirement)}}\\ \hline
TC-07 & Chủ sở hữu chuyển nhượng NFT lô tín chỉ & YC-?? & Đạt\\ \hline
TC-08 & Hủy một phần lô, cấp chứng nhận với dãy sê-ri đúng & YC-?? & Đạt\\ \hline
TC-09 & Hủy hết lô thì NFT bị đốt, không chuyển nhượng hay hủy lại được (sử dụng kép) & YC-?? & Đạt\\ \hline
TC-10 & Không hủy quá số còn lại, không hủy 0, bắt buộc có người thụ hưởng & YC-?? & Đạt\\ \hline
TC-11 & Người không sở hữu không hủy được; người được approve thì hủy được & YC-?? & Đạt\\ \hline
TC-12 & retireFor chỉ dành cho RETIREMENT\_AGENT\_ROLE đang giữ NFT & YC-?? & Đạt\\ \hline
TC-13 & getBatch của lô không tồn tại bị từ chối & YC-?? & Đạt\\ \hline
TC-14 & Cơ quan quản lý tạm dừng thì chặn mint, chuyển nhượng và hủy & YC-?? & Đạt\\ \hline
TC-15 & Hỗ trợ các interface ERC-721 và AccessControl & YC-?? & Đạt\\ \hline
\multicolumn{4}{|l|}{\textit{Pool ERC-20 (giao dịch từng phần)}}\\ \hline
TC-16 & Nạp NFT vào pool nhận đúng số token, NFT bị khóa trong pool & YC-?? & Đạt\\ \hline
TC-17 & Nạp bằng safeTransferFrom trực tiếp; lô đã hủy một phần chỉ nhận phần còn lại & YC-?? & Đạt\\ \hline
TC-18 & Pool từ chối NFT không phải của hợp đồng tín chỉ & YC-?? & Đạt\\ \hline
TC-19 & Chuyển token ERC-20 từng phần cho người mua & YC-?? & Đạt\\ \hline
TC-20 & Người mua hủy từ pool, chứng nhận ghi người mua và dãy sê-ri; token đã đốt không dùng lại được (sử dụng kép) & YC-?? & Đạt\\ \hline
TC-21 & Bất biến: tổng cung token bằng tổng phần còn lại của các lô trong pool & YC-?? & Đạt\\ \hline
TC-22 & Đổi lại (redeem) lô cần đốt đủ token bằng phần còn lại & YC-?? & Đạt\\ \hline
TC-23 & Thao tác với lô không nằm trong pool bị từ chối & YC-?? & Đạt\\ \hline
TC-24 & NFT trong pool không thể bị chủ cũ bán riêng (bán kép) & YC-?? & Đạt\\ \hline
TC-25 & Hủy từ pool quá phần còn lại của lô bị từ chối, số dư không đổi & YC-?? & Đạt\\ \hline
TC-26 & Nạp lô đã hủy hết (không còn tồn tại) bị từ chối & YC-?? & Đạt\\ \hline
\end{longtable}
\endgroup}

\newcommand{\BangTestAPI}{%
\begingroup\small
\renewcommand{\arraystretch}{1.12}
\begin{longtable}{|p{1.35cm}|p{8.7cm}|p{1.55cm}|p{1.55cm}|}
\caption{Kết quả kiểm thử API và bộ đồng bộ sự kiện}\label{tab:testapi}\\
\hline
\textbf{Mã} & \textbf{Kịch bản kiểm thử} & \textbf{Yêu cầu} & \textbf{Kết quả}\\ \hline
\endfirsthead
\hline
\textbf{Mã} & \textbf{Kịch bản kiểm thử} & \textbf{Yêu cầu} & \textbf{Kết quả}\\ \hline
\endhead
API-01 & Đăng nhập bằng chữ ký ví, tài khoản mới tự chọn vai trò & YC-?? & Đạt\\ \hline
API-02 & Chữ ký sai hoặc dùng lại nonce bị từ chối & YC-?? & Đạt\\ \hline
API-03 & Không tự đăng ký được vai trò VERIFIER; ADMIN cấp được & YC-?? & Đạt\\ \hline
API-04 & Endpoint quản trị từ chối người không phải ADMIN và người chưa đăng nhập & YC-?? & Đạt\\ \hline
API-05 & Chủ dự án đăng ký dự án; trùng mã dự án bị từ chối & YC-?? & Đạt\\ \hline
API-06 & Gửi yêu cầu cấp tín chỉ cần báo cáo giám sát đã tải lên & YC-?? & Đạt\\ \hline
API-07 & Chỉ VVB duyệt được; duyệt xong trả về tham số mint với sê-ri đề xuất & YC-?? & Đạt\\ \hline
API-08 & Báo cáo thẩm định đã dùng không duyệt được yêu cầu khác (phát hành kép, off-chain) & YC-?? & Đạt\\ \hline
API-09 & Sự kiện CreditMinted tạo đợt cấp, khối tín chỉ và chuyển yêu cầu sang MINTED & YC-?? & Đạt\\ \hline
API-10 & Sự kiện trùng (cùng tx\_hash, log\_index) không bị ghi hai lần & YC-?? & Đạt\\ \hline
API-11 & Sự kiện CreditRetired tách khối và tạo chứng nhận tra cứu được & YC-?? & Đạt\\ \hline
API-12 & Tra cứu theo số sê-ri cho biết sê-ri đã hủy hay còn hiệu lực & YC-?? & Đạt\\ \hline
API-13 & Dashboard so sánh tín chỉ đã hủy với phần phát thải vượt hạn ngạch & YC-?? & Đạt\\ \hline
API-14 & Hủy hết phần còn lại thì không còn khối Held; thống kê khớp & YC-?? & Đạt\\ \hline
API-15 & Xuất dữ liệu dự án theo cấu trúc CAD Trust & YC-?? & Đạt\\ \hline
API-16 & Metadata NFT và tài liệu công khai cho phép đối chiếu mã băm & YC-?? & Đạt\\ \hline
E2E-01 & Bộ đồng bộ đọc sự kiện thật trên mạng Hardhat (mint, nạp pool, hủy); chạy lại không ghi trùng; phần còn lại on-chain khớp CSDL & YC-?? & Đạt\\ \hline
E2E-02 & Kiểm thử toàn hệ thống đang chạy: đăng nhập ví, đăng ký dự án, yêu cầu, VVB duyệt và mint, nạp pool, bán, hủy, tra cứu, dashboard & YC-?? & Đạt\\ \hline
\end{longtable}
\endgroup}

\newcommand{\BangDoPhu}{%
\begin{table}[htbp]
\centering
\caption{Độ phủ mã của bộ kiểm thử hợp đồng (solidity-coverage)}
\label{tab:dophu}
\begin{tabular}{|l|r|r|r|r|}
\hline
\textbf{Tệp} & \textbf{Câu lệnh} & \textbf{Nhánh} & \textbf{Hàm} & \textbf{Dòng}\\ \hline
CarbonCreditNFT.sol & 100\% & 94,44\% & 100\% & 100\%\\ \hline
CarbonPool.sol & 100\% & 75,00\% & 100\% & 100\%\\ \hline
\textbf{Tổng} & \textbf{100\%} & \textbf{88,46\%} & \textbf{100\%} & \textbf{100\%}\\ \hline
\end{tabular}
\end{table}}

\newcommand{\BangGas}{%
\begin{table}[htbp]
\centering
\caption{Chi phí gas của các thao tác chính (Solidity 0.8.24, optimizer 200 runs)}
\label{tab:gas}
\small
\begin{tabular}{|l|l|r|r|r|}
\hline
\textbf{Hợp đồng} & \textbf{Thao tác} & \textbf{Nhỏ nhất} & \textbf{Lớn nhất} & \textbf{Trung bình}\\ \hline
CarbonCreditNFT & Triển khai & -- & -- & 2.431.505\\ \hline
CarbonPool & Triển khai & -- & -- & 1.167.804\\ \hline
CarbonCreditNFT & mintBatch & 260.483 & 294.683 & 288.105\\ \hline
CarbonCreditNFT & retire (trực tiếp) & 206.215 & 261.001 & 243.049\\ \hline
CarbonCreditNFT & transferFrom & -- & -- & 57.310\\ \hline
CarbonCreditNFT & safeTransferFrom vào pool & -- & -- & 182.071\\ \hline
CarbonPool & deposit & 116.482 & 184.882 & 157.522\\ \hline
CarbonPool & transfer (PCC) & 29.683 & 51.595 & 47.208\\ \hline
CarbonPool & retire (từ pool) & 262.346 & 282.681 & 274.041\\ \hline
CarbonPool & redeem & -- & -- & 84.162\\ \hline
\end{tabular}
\end{table}}

\newcommand{\BangTinhToanKep}{%
\begin{table}[htbp]
\centering
\caption{Kịch bản chống tính toán kép và cơ chế chặn}
\label{tab:tinhtoankep}
\small
\renewcommand{\arraystretch}{1.2}
\begin{tabularx}{\textwidth}{|p{2.6cm}|X|X|p{1.9cm}|}
\hline
\textbf{Dạng} & \textbf{Kịch bản tấn công} & \textbf{Cơ chế chặn} & \textbf{Test}\\ \hline
Phát hành kép & VVB mint lô có dãy sê-ri chồng lên lô đã cấp của cùng dự án và vintage & nextAvailableSerial: sê-ri phải tăng dần, lỗi SerialOverlap & TC-03\\ \hline
Phát hành kép & Dùng lại một báo cáo thẩm định để cấp lần hai (kể cả cho dự án khác) & documentUsed on-chain (DocumentAlreadyUsed) và verifications.document\_hash unique off-chain & TC-05, API-08\\ \hline
Sử dụng kép & Hủy lại hoặc chuyển nhượng tín chỉ đã hủy & Trừ phần đã hủy khỏi lô; hủy hết thì đốt NFT & TC-09, TC-10\\ \hline
Sử dụng kép & Dùng lại token PCC đã đốt khi hủy từ pool & \_burn trước khi gọi retireFor; số dư không đủ thì giao dịch bị hoàn tác & TC-20\\ \hline
Bán kép & Chủ cũ bán riêng NFT đã nạp vào pool (đã nhận token tương ứng) & NFT thuộc sở hữu của pool, chỉ ra qua redeem khi đốt đủ token & TC-24, TC-21\\ \hline
Tuyên bố kép (trong hệ thống) & Hai bên cùng nhận đã bù trừ bằng một dãy sê-ri & Mỗi chứng nhận ghi dãy sê-ri và người thụ hưởng duy nhất; tra cứu công khai theo sê-ri & TC-08, API-12\\ \hline
\end{tabularx}
\end{table}}


% Tổng quan (Normal, không tiêu đề): chương trình bày các công nghệ đồ án dùng để hiện thực thiết kế ở
% Chương 3, kết quả cài đặt, kiểm thử và triển khai trên Sepolia.
% [VIẾT SAU] \canviet{Đoạn tổng quan chương}

% ============ PHẦN A — CÔNG NGHỆ SỬ DỤNG (từ Chương 2 cũ) ============
% [VIẾT SAU] Toàn bộ các mục dưới đây chưa viết nội dung — bỏ dấu % để khôi phục.
% % ------------------------------------------------------------
% \section{Blockchain và nền tảng Ethereum}
% \label{sec:blockchain}
% \label{sec:ethereum}
%
% \subsection*{Blockchain}
% \label{subsec:cautrucblock}
% \label{subsec:dongthuan}
% % Ý cần có (ngắn gọn, chỉ phần đồ án dùng): khối, mã băm khối trước, sổ cái chỉ ghi thêm (ledger) ->
% % vì sao dữ liệu on-chain khó sửa, phục vụ yêu cầu toàn vẹn ở mục \ref{sec:phichucnang} và VĐ-2;
% % đồng thuận PoW/PoS, Ethereum chuyển sang PoS (The Merge 15/9/2022).
% % Trích: \cite{nakamoto2008bitcoin,zheng2017overview,ethereum2024pos}. Nhắc \autoref{fig:chuoikhoi} trước hình.
% % [VIẾT SAU] \canviet{Blockchain}
% %
% % \HinhChuoiKhoi
%
% \subsection*{Lựa chọn nền tảng blockchain}
% \label{subsec:phanloai}
% % Bảng so sánh Ethereum (public) với Hyperledger Fabric (permissioned, mà Vaccargiu dùng) và Chia
% % (CAD Trust dùng): mô hình truy cập, đồng thuận, ngôn ngữ smart contract, chuẩn token, ví, testnet, chi phí.
% % Lý do chọn Ethereum: phiếu giao nhiệm vụ; chuẩn ERC-721/ERC-20; MetaMask; testnet công khai để bên thứ ba
% % tự kiểm tra (VĐ-2, NC-07). Nhược điểm: phí gas, dữ liệu công khai -> xử lý ở yêu cầu quyền riêng tư.
% % [VIẾT SAU] \canviet{Lựa chọn nền tảng blockchain}
%
% \subsection*{Tài khoản, giao dịch và phí gas}
% \label{subsec:taikhoan}
% \label{subsec:evm}
% % Ý cần có: EOA và contract account; khóa riêng, địa chỉ ví; giao dịch ký số, nonce; EVM, gas, EIP-1559.
% % Đây là cơ sở cho phần đo gas ở Chương 3. Trích: \cite{buterin2014ethereum,wood2014yellow,eip1559}.
% % [VIẾT SAU] \canviet{Tài khoản, giao dịch và phí gas}
%
% \subsection*{Testnet Sepolia và Etherscan}
% \label{subsec:sepolia}
% % Ý cần có: mainnet và testnet; Sepolia, faucet; Etherscan để tra giao dịch, xác minh (verify) mã hợp đồng.
% % Lý do: phạm vi nguyên mẫu (mục \ref{sec:muctieu}), không tốn phí thật.
% % [VIẾT SAU] \canviet{Testnet Sepolia}
%
% % ------------------------------------------------------------
% \section{Smart contract và công cụ phát triển}
% \label{sec:smartcontract}
%
% \subsection*{Smart contract và ngôn ngữ Solidity}
% \label{subsec:khainiemsc}
% \label{subsec:solidity}
% % Ý cần có: định nghĩa (Szabo); tự thực thi, không sửa được mã sau triển khai; event và log (dùng cho
% % bộ đồng bộ sự kiện); cấu trúc contract (state variable, function, modifier, event), kế thừa.
% % Yêu cầu phục vụ: NC-03, NC-06 (ràng buộc sê-ri, không hủy hai lần thành ràng buộc kỹ thuật).
% % Trích: \cite{szabo1997formalizing,buterin2014ethereum,solidity2024docs}.
% % [VIẾT SAU] \canviet{Smart contract và Solidity}
%
% \subsection*{Thư viện OpenZeppelin và Hardhat}
% \label{subsec:openzeppelin}
% % OpenZeppelin 5: ERC721, ERC20, AccessControl (VERIFIER_ROLE...), Pausable, ReentrancyGuard -> NC-09 và
% % yêu cầu bảo mật. Hardhat: biên dịch, Hardhat Network, test Mocha/Chai, coverage, gas reporter.
% % Phương án thay thế: Foundry, Truffle (đã ngừng phát triển); lý do chọn Hardhat (JavaScript cùng ngôn ngữ backend).
% % Trích: \cite{openzeppelin2024contracts}.
% % [VIẾT SAU] \canviet{OpenZeppelin và Hardhat}
%
% \subsection*{Các lỗ hổng thường gặp của smart contract}
% \label{subsec:baomat}
% % Reentrancy, kiểm soát truy cập sai, tràn số (Solidity >= 0.8 tự kiểm tra), tx.origin. Mỗi lỗ hổng
% % 1–2 câu và cách phòng trong đồ án; công cụ quét Slither. Trích: \cite{atzei2017survey}.
% % [VIẾT SAU] \canviet{Lỗ hổng smart contract}
%
% % ------------------------------------------------------------
% \section{Các chuẩn token}
% \label{sec:token-standards}
%
% \subsection*{ERC-20 và ERC-721}
% \label{subsec:erc20}
% \label{subsec:erc721}
% % ERC-20: fungible token, totalSupply, balanceOf, transfer, approve, transferFrom, allowance \cite{eip20}.
% % ERC-721: NFT, ownerOf, safeTransferFrom, tokenURI, onERC721Received \cite{eip721}.
% % [VIẾT SAU] \canviet{ERC-20 và ERC-721}
%
% \subsection*{Lựa chọn chuẩn token cho tín chỉ các-bon}
% \label{subsec:chontoken}
% % Bảng so sánh ERC-20 / ERC-721 / ERC-1155 theo: giữ metadata riêng từng lô, chia nhỏ, thanh khoản,
% % hỗ trợ ví. Kết luận: mỗi lô là một NFT ERC-721 (giữ dự án, vintage, dãy sê-ri, NC-03), pool ERC-20
% % cho giao dịch từng phần (NC-05); kế thừa mô hình Batch NFT/TCO2/BCT của Toucan (mục \ref{sec:khaosatsobo}).
% % [VIẾT SAU] \canviet{Lựa chọn chuẩn token}
%
% % ------------------------------------------------------------
% \section{Ứng dụng web giao tiếp với blockchain}
% \label{sec:dapp}
%
% \subsection*{Kiến trúc dApp}
% \label{subsec:kientrucdapp}
% % Luồng một giao dịch từ giao diện tới smart contract (giải thích \autoref{fig:luongdapp}).
% % [VIẾT SAU] \canviet{Kiến trúc dApp}
% %
% % \HinhLuongDApp
%
% \subsection*{Ví MetaMask, ethers.js và Next.js}
% \label{subsec:metamask}
% \label{subsec:ethers}
% % MetaMask giữ khóa riêng, EIP-1193 (window.ethereum) \cite{eip1193}; ethers.js: Provider, Signer,
% % Contract, ABI, call và send \cite{ethers2024docs}; Next.js cho giao diện. Yêu cầu phục vụ: tính dễ dùng,
% % người dùng tự ký giao dịch (backend không giữ khóa). Phương án thay thế: web3.js, viem; React thuần.
% % [VIẾT SAU] \canviet{MetaMask, ethers.js và Next.js}
%
% % ------------------------------------------------------------
% \section{Công nghệ phía off-chain}
% \label{sec:offchain-congnghe}
% \label{sec:onoffchain-lythuyet}
% % Node.js + Express (API RESTful), PostgreSQL (dữ liệu theo CAD Trust V2.0, NC-10), JWT, xác thực bằng
% % chữ ký ví. Tiêu chí chọn dữ liệu on-chain/off-chain theo \cite{xu2017taxonomy}; neo mã băm keccak256 của hồ sơ;
% % IPFS \cite{benet2014ipfs} là hướng mở rộng (hiện lưu tệp cục bộ). Phương án thay thế: MongoDB; lý do chọn
% % CSDL quan hệ (mô hình CAD Trust có quan hệ rõ, giống UNDP dùng PostgreSQL).
% % [VIẾT SAU] \canviet{Công nghệ off-chain}


% ============ PHẦN B — XÂY DỰNG, KIỂM THỬ, TRIỂN KHAI (từ Chương 3 cũ) ============
% ------------------------------------------------------------
\section{Xây dựng ứng dụng}
\label{sec:xaydung}

\subsection*{Thư viện và công cụ sử dụng}
\label{sec:moitruong}
% Mẫu SoICT: bảng công cụ có phiên bản (và URL). Giải thích \autoref{tab:congnghe}; lý do chọn công cụ đã nêu
% ở Chương 2 nên chỉ nhắc lại ngắn. Cấu trúc mã nguồn theo \autoref{fig:cautrucma}.
% [VIẾT SAU] \canviet{Thư viện và công cụ sử dụng}

\BangCongNghe

\HinhCauTrucMa

\subsection*{Kết quả đạt được}
\label{subsec:ketquadatduoc}
\label{sec:caidat-hopdong}
\label{sec:caidat-backend}
% Mẫu SoICT: mô tả các sản phẩm đóng gói, thành phần, vai trò; bảng thống kê số dòng mã, số tệp, số hợp đồng,
% số endpoint, dung lượng.
% Nội dung cài đặt từ Chương 4 cũ (chỉ tóm lược, chi tiết đóng góp ở Chương 4):
%  - CarbonCreditNFT: struct CreditBatch (projectId, vintage, serialFrom, quantity, retired, documentHash,
%    verifier, issuedAt); mintBatch: seriesKey = keccak256(projectId, vintage), serialFrom >= nextAvailableSerial,
%    documentUsed[documentHash]; retire hủy từ đầu dãy sê-ri còn lại; retireFor cho pool (RETIREMENT_AGENT_ROLE);
%    Pausable ghi đè _update.
%  - CarbonPool: decimals = 0; nạp qua onERC721Received; redeem đốt remainingOf(tokenId); retire đốt token rồi
%    gọi retireFor; ReentrancyGuard.
%  - Phân quyền: DEFAULT_ADMIN_ROLE và vai trò phát hành ISSUER_ROLE (cơ quan quản lý, mỗi tỉnh một tài khoản),
%    VERIFIER_ROLE (VVB, chỉ phê duyệt off-chain), RETIREMENT_AGENT_ROLE (pool). [CẦN SỬA MÃ] thêm ISSUER_ROLE.
%  - Backend: db/schema.sql từ ERD; xác thực nonce -> personal_sign -> ethers.verifyMessage -> JWT 8 giờ;
%    API theo \autoref{tab:api} (backend không giữ khóa riêng); indexer quét queryFilter, khóa (tx_hash, log_index),
%    chưa xử lý reorg.
% [VIẾT SAU] \canviet{Kết quả đạt được}

% \BangAPI   % [17/11/2026] đã dùng ở Chương 3, mục Thiết kế API

\subsection*{Minh họa các chức năng chính}
\label{subsec:minhhoa}
\label{sec:caidat-frontend}
% Mỗi màn hình 1 ảnh chụp + giải thích ngắn: danh sách dự án, chi tiết dự án, tra cứu sê-ri, chứng nhận hủy,
% dự án của tôi, thẩm định (VVB), ví tín chỉ, dashboard, quản trị. Chụp khi chạy thật (README "Chạy thử cục bộ").
% [VIẾT SAU] \canviet{Minh họa các chức năng chính}

% ------------------------------------------------------------
\section{Kiểm thử}
\label{sec:kiemthu}
\label{sec:pp-kiemthu}
\label{sec:kq-chucnang}
% Mẫu SoICT: 2–3 trang; thiết kế ca kiểm thử cho 2–3 chức năng quan trọng nhất (phát hành, hủy từ pool, tra cứu),
% nêu kỹ thuật kiểm thử; tổng kết số ca và kết quả; giải thích ca không đạt. Ca kiểm thử khác đưa vào Phụ lục.
% Ba mức: kiểm thử đơn vị hợp đồng trên Hardhat Network (Mocha/Chai, loadFixture); kiểm thử API với pg-mem và
% supertest; kiểm thử tích hợp e2e-indexer và smoke-api. Ghi cấu hình máy, phiên bản Node, Hardhat.
% Nhận xét độ phủ: nhánh của CarbonPool 75% vì ... (xem coverage/index.html).
% Chi phí gas: quy đổi ETH/USD theo giá gas tại thời điểm đo (ghi nguồn, ngày); vì sao mintBatch, retire tốn gas.
% Kịch bản chống tính toán kép (\BangTinhToanKep) đặt ở Chương 4, mục \ref{sec:gp-tinhtoankep}.
% [VIẾT SAU] \canviet{Kiểm thử}

\BangTestHopDong

\BangTestAPI

\BangDoPhu

\BangGas

% ------------------------------------------------------------
\section{Triển khai}
\label{sec:trienkhai}
\label{sec:kq-sepolia}
% Giải thích \autoref{fig:trienkhai}: Sepolia + RPC, backend + PostgreSQL, frontend; địa chỉ hợp đồng,
% link Etherscan (đã verify), vài giao dịch mẫu (mint, retire), ảnh chụp.
% [VIẾT SAU] \canviet{Triển khai}

\HinhTrienKhai


% ------------------------------------------------------------
% Kết chương (Normal, không tiêu đề): tóm tắt bộ công nghệ (bảng Công nghệ | Phiên bản | Vai trò | Yêu cầu
% phục vụ), kết quả cài đặt, kiểm thử, triển khai; dẫn sang Chương 5 (giải pháp và đóng góp).
% [VIẾT SAU] \canviet{Kết chương 4}
 để khôi phục.
% Bảng API (\BangAPI) đã dùng ở Chương 3 (Thiết kế API), không đặt lại ở đây.
%
% Hướng dẫn của mẫu SoICT cho phần công nghệ: không quá 10 trang; chỉ giới thiệu công nghệ đồ án thực sự
% dùng; với mỗi công nghệ nêu (i) yêu cầu nó phục vụ (NC-xx, VĐ-x, yêu cầu phi chức năng ở Chương 2),
% (ii) phương án thay thế, (iii) lý do chọn.
% Văn phong: xem ghi chú đầu Chap1.tex. [VIẾT SAU] toàn chương.
% ============================================================

\chapter{CÔNG NGHỆ, XÂY DỰNG VÀ KIỂM THỬ HỆ THỐNG}
\chaptermark{CÔNG NGHỆ, XÂY DỰNG VÀ KIỂM THỬ}
\label{chap:congnghe}
\label{chap:cosolythuyet}
\label{chap:xaydung}
\label{chap:ketqua}
% ============================================================
% HÌNH VÀ BẢNG CHƯƠNG 4 (CÔNG NGHỆ, XÂY DỰNG VÀ KIỂM THỬ) — gọi bằng macro trong Chap4.tex
%   \HinhChuoiKhoi  : cấu trúc liên kết các khối           (fig:chuoikhoi)   [chuyển từ figures_chap2 ngày 17/11/2026]
%   \HinhLuongDApp  : luồng tương tác của dApp              (fig:luongdapp)   [chuyển từ figures_chap2 ngày 17/11/2026]
%   \BangCongNghe   : công nghệ và phiên bản sử dụng      (tab:congnghe)
%   \HinhTrienKhai  : sơ đồ triển khai hệ thống            (fig:trienkhai)
%   \HinhCauTrucMa  : cấu trúc thư mục mã nguồn            (fig:cautrucma)
% Phiên bản lấy từ package.json / hardhat.config.js tại thời điểm 21/10/2026.
% ============================================================

% ---------- Cấu trúc liên kết các khối ----------
\newcommand{\HinhChuoiKhoi}{%
\begin{figure}[htbp]
\centering
\begin{tikzpicture}[
    >=Stealth,
    truong/.style={draw, minimum width=3.3cm, minimum height=0.55cm,
                   font=\footnotesize, fill=white, inner sep=1pt},
    tieude/.style={truong, fill=gray!20, font=\footnotesize\bfseries}]
  \foreach \i/\x/\ten in {0/0/N-1, 1/4.6/N, 2/9.2/N+1} {
    \node[tieude] (h\i) at (\x,0) {Khối \ten};
    \node[truong, below=0pt of h\i] (p\i) {Mã băm khối trước};
    \node[truong, below=0pt of p\i] (t\i) {Dấu thời gian};
    \node[truong, below=0pt of t\i] (m\i) {Gốc Merkle};
    \node[truong, below=0pt of m\i, minimum height=1cm, align=center] (g\i)
        {Danh sách\\giao dịch};
    \node[draw, thick, rounded corners=2pt, inner sep=3pt, fit=(h\i)(g\i)] (b\i) {};
  }
  \draw[->, thick] (p1.west) -- (b0.east |- p1.west);
  \draw[->, thick] (p2.west) -- (b1.east |- p2.west);
  \node[font=\large] at (-2.35,-1.4) {$\cdots$};
  \node[font=\large] at (11.55,-1.4) {$\cdots$};
\end{tikzpicture}
\caption[Cấu trúc liên kết giữa các khối trong chuỗi khối]{Cấu trúc liên kết giữa các khối trong chuỗi khối \hinhtao[Claude]{20/10/2026}}
\label{fig:chuoikhoi}
\end{figure}}

% ---------- Luồng tương tác của ứng dụng phi tập trung ----------
\newcommand{\HinhLuongDApp}{%
\begin{figure}[htbp]
\centering
\begin{tikzpicture}[
    >=Stealth,
    hop/.style={draw, rounded corners=3pt, align=center, font=\footnotesize,
                text width=2.1cm, minimum height=1.1cm, fill=white},
    nhan/.style={font=\scriptsize, align=center}]
  \node[hop]                         (u)   at (0,0)     {Người dùng};
  \node[hop]                         (fe)  at (3.7,0)   {Giao diện web\\(Next.js,\\ethers.js)};
  \node[hop]                         (mm)  at (7.4,0)   {Ví\\MetaMask\\(giữ khóa riêng)};
  \node[hop, text width=2.6cm, minimum height=1.6cm, fill=gray!12]
                                     (net) at (11.3,0)  {Mạng Ethereum\\(Sepolia)\\nút RPC, EVM,\\hợp đồng thông minh};
  \node[hop]                         (be)  at (3.7,-3)  {Backend\\(Node.js, Express)};
  \node[hop]                         (db)  at (0,-3)    {CSDL off-chain\\(PostgreSQL)};

  \draw[->] (u) -- node[nhan, above] {thao tác} (fe);
  \draw[->] (fe) -- node[nhan, above] {yêu cầu ký\\(EIP-1193)} (mm);
  \draw[->] (mm) -- node[nhan, above] {giao dịch\\đã ký} (net);
  \draw[->, dashed] (fe.north) to[bend left=22]
      node[nhan, above] {đọc trạng thái on-chain (call)} (net.north);
  \draw[->] (net.south) |- node[nhan, pos=0.75, below] {sự kiện (Minted, Transfer, Retired\ldots)} (be.east);
  \draw[<->] (fe) -- node[nhan, right] {API REST} (be);
  \draw[<->] (be) -- node[nhan, above] {đọc/ghi} (db);
\end{tikzpicture}
\caption[Luồng tương tác giữa giao diện, ví MetaMask, mạng Ethereum và backend]{Luồng tương tác giữa giao diện, ví MetaMask, mạng Ethereum và backend \hinhtao[Claude]{20/10/2026}}
\label{fig:luongdapp}
\end{figure}}

\newcommand{\BangCongNghe}{%
\begin{table}[htbp]
\centering
\caption{Công nghệ và phiên bản sử dụng}
\label{tab:congnghe}
\small
\renewcommand{\arraystretch}{1.2}
\begin{tabularx}{\textwidth}{|p{3.2cm}|p{4.4cm}|X|}
\hline
\textbf{Thành phần} & \textbf{Công nghệ (phiên bản)} & \textbf{Vai trò}\\ \hline
Hợp đồng thông minh & Solidity 0.8.24, OpenZeppelin Contracts 5.1.0 & ERC-721, ERC-20, AccessControl, Pausable, ReentrancyGuard\\ \hline
Môi trường phát triển & Hardhat 2.29, hardhat-toolbox 5 (ethers 6, Mocha/Chai, solidity-coverage, gas reporter) & Biên dịch, kiểm thử, đo độ phủ và gas, triển khai\\ \hline
Mạng blockchain & Hardhat Network (cục bộ), Ethereum Sepolia & Kiểm thử và triển khai thử nghiệm\\ \hline
Backend & Node.js 22, Express 4.21, jsonwebtoken 9, multer & API REST, xác thực bằng chữ ký ví, tải hồ sơ\\ \hline
Cơ sở dữ liệu & PostgreSQL 16 (pg 8.13); pg-mem khi kiểm thử & Lưu dữ liệu off-chain theo mô hình CAD Trust\\ \hline
Đồng bộ sự kiện & ethers 6.13 (JsonRpcProvider, queryFilter) & Đọc sự kiện on-chain vào CSDL\\ \hline
Frontend & Next.js 14.2 (App Router), React 18.3, ethers 6.13 & Giao diện web, kết nối MetaMask\\ \hline
Ví & MetaMask & Ký đăng nhập và ký giao dịch\\ \hline
\end{tabularx}
\end{table}}

\newcommand{\HinhTrienKhai}{%
\begin{figure}[htbp]
\centering
\begin{adjustbox}{max width=\textwidth}
\begin{tikzpicture}[>=Stealth, font=\footnotesize,
    hop/.style={draw, rounded corners=3pt, align=center, fill=white, minimum height=1.1cm, text width=3.2cm},
    nhan/.style={font=\scriptsize, align=center}]
  \node[hop] (br)  at (0,0)    {\textbf{Trình duyệt}\\ giao diện Next.js\\ + MetaMask};
  \node[hop] (fe)  at (0,-3.2) {\textbf{Máy chủ frontend}\\ Next.js (Vercel)};
  \node[hop] (be)  at (5.2,0)  {\textbf{Máy chủ backend}\\ Express, port 4000\\ + bộ đồng bộ sự kiện};
  \node[hop] (db)  at (10.4,0) {\textbf{PostgreSQL}};
  \node[hop] (fs)  at (10.4,-3.2) {\textbf{Kho tệp}\\ uploads/ (hoặc IPFS)};
  \node[hop] (rpc) at (5.2,-3.2) {\textbf{Nút RPC}\\ (Infura/Alchemy)};
  \node[hop, fill=gray!10] (eth) at (5.2,-6.0) {\textbf{Ethereum Sepolia}\\ CarbonCreditNFT, CarbonPool};
  \draw[->] (fe) -- node[nhan, left] {tải trang\\ (HTTPS)} (br);
  \draw[<->] (br) -- node[nhan, above] {REST/JSON\\ + JWT} (be);
  \draw[<->] (be) -- node[nhan, above] {SQL} (db);
  \draw[->] (be) -- node[nhan, right] {lưu tệp} (fs);
  \draw[->] (be) -- node[nhan, right] {đọc sự kiện} (rpc);
  \draw[->] (br.south east) to[bend left=10] node[nhan, pos=0.35, right] {MetaMask ký\\ và gửi giao dịch} (rpc.north west);
  \draw[<->] (rpc) -- (eth);
\end{tikzpicture}
\end{adjustbox}
\caption[Sơ đồ triển khai hệ thống]{Sơ đồ triển khai hệ thống \hinhtao[Claude]{21/10/2026}}
\label{fig:trienkhai}
\end{figure}}

\newcommand{\HinhCauTrucMa}{%
\begin{figure}[htbp]
\centering
\footnotesize
\renewcommand{\arraystretch}{1.05}
\begin{tabular}{@{}l l@{}}
\texttt{DATN-CarbonCreditOnBlockChain/} & \\
\texttt{\quad contracts/} & \textit{Hardhat}\\
\texttt{\quad\quad contracts/CarbonCreditNFT.sol} & lô tín chỉ ERC-721: mint, retire\\
\texttt{\quad\quad contracts/CarbonPool.sol} & pool ERC-20: deposit, redeem, retire\\
\texttt{\quad\quad test/CarbonCredit.test.js} & 26 test case TC-01 đến TC-26\\
\texttt{\quad\quad scripts/deploy.js} & triển khai, cấp quyền\\
\texttt{\quad\quad scripts/demo.js} & kịch bản demo, đo gas\\
\texttt{\quad\quad scripts/e2e-indexer.js} & kiểm thử on-chain $\rightarrow$ CSDL\\
\texttt{\quad\quad scripts/smoke-api.js} & kiểm thử toàn hệ thống đang chạy\\
\texttt{\quad backend/} & \textit{Node.js + Express}\\
\texttt{\quad\quad db/schema.sql} & lược đồ CSDL (CAD Trust + ứng dụng)\\
\texttt{\quad\quad src/routes/*.js} & các nhóm endpoint API\\
\texttt{\quad\quad src/indexer.js} & đồng bộ sự kiện on-chain\\
\texttt{\quad\quad test/api.test.js} & 16 test case API-01 đến API-16\\
\texttt{\quad frontend/} & \textit{Next.js}\\
\texttt{\quad\quad app/} & các trang giao diện\\
\texttt{\quad\quad components/WalletProvider.js} & đăng nhập bằng ví\\
\texttt{\quad\quad lib/chain.js, lib/api.js} & gọi hợp đồng và API\\
\end{tabular}
\caption[Cấu trúc thư mục mã nguồn]{Cấu trúc thư mục mã nguồn \hinhtao[Claude]{21/10/2026}}
\label{fig:cautrucma}
\end{figure}}

% ============================================================
% BẢNG CHƯƠNG 5 — số liệu lấy trực tiếp từ kết quả chạy ngày 21/10/2026
%   (npx hardhat test / coverage / REPORT_GAS, npm test trong backend)
%   \BangTestHopDong : 26 test case hợp đồng               (tab:testhopdong)
%   \BangTestAPI     : 16 test case API                    (tab:testapi)
%   \BangDoPhu       : độ phủ mã hợp đồng                  (tab:dophu)
%   \BangGas         : chi phí gas các thao tác chính      (tab:gas)
%   \BangTinhToanKep : kịch bản chống tính toán kép        (tab:tinhtoankep)
% CHẠY LẠI và cập nhật số liệu nếu sửa hợp đồng. Cột "Yêu cầu" cần điền mã YC-xx
% sau khi bạn lập bảng yêu cầu chức năng ở mục \ref{subsec:ycchucnang}.
% ============================================================

\newcommand{\BangTestHopDong}{%
\begingroup\small
\renewcommand{\arraystretch}{1.12}
\begin{longtable}{|p{1.25cm}|p{8.8cm}|p{1.55cm}|p{1.55cm}|}
\caption{Kết quả kiểm thử đơn vị hợp đồng thông minh}\label{tab:testhopdong}\\
\hline
\textbf{Mã} & \textbf{Kịch bản kiểm thử} & \textbf{Yêu cầu} & \textbf{Kết quả}\\ \hline
\endfirsthead
\hline
\textbf{Mã} & \textbf{Kịch bản kiểm thử} & \textbf{Yêu cầu} & \textbf{Kết quả}\\ \hline
\endhead
\multicolumn{4}{|l|}{\textit{Phát hành (Minting)}}\\ \hline
TC-01 & VVB mint lô tín chỉ, lưu đúng thông tin và phát sự kiện CreditMinted & YC-?? & Đạt\\ \hline
TC-02 & Tài khoản không có VERIFIER\_ROLE không mint được & YC-?? & Đạt\\ \hline
TC-03 & Chặn dãy sê-ri chồng lấn trong cùng dự án và vintage (phát hành kép) & YC-?? & Đạt\\ \hline
TC-04 & Cho phép lô kế tiếp không chồng lấn, và cùng dãy sê-ri ở vintage khác & YC-?? & Đạt\\ \hline
TC-05 & Chặn dùng lại cùng một báo cáo thẩm định (phát hành kép) & YC-?? & Đạt\\ \hline
TC-06 & Từ chối số lượng 0 và sê-ri 0 & YC-?? & Đạt\\ \hline
\multicolumn{4}{|l|}{\textit{Chuyển nhượng và hủy (Retirement)}}\\ \hline
TC-07 & Chủ sở hữu chuyển nhượng NFT lô tín chỉ & YC-?? & Đạt\\ \hline
TC-08 & Hủy một phần lô, cấp chứng nhận với dãy sê-ri đúng & YC-?? & Đạt\\ \hline
TC-09 & Hủy hết lô thì NFT bị đốt, không chuyển nhượng hay hủy lại được (sử dụng kép) & YC-?? & Đạt\\ \hline
TC-10 & Không hủy quá số còn lại, không hủy 0, bắt buộc có người thụ hưởng & YC-?? & Đạt\\ \hline
TC-11 & Người không sở hữu không hủy được; người được approve thì hủy được & YC-?? & Đạt\\ \hline
TC-12 & retireFor chỉ dành cho RETIREMENT\_AGENT\_ROLE đang giữ NFT & YC-?? & Đạt\\ \hline
TC-13 & getBatch của lô không tồn tại bị từ chối & YC-?? & Đạt\\ \hline
TC-14 & Cơ quan quản lý tạm dừng thì chặn mint, chuyển nhượng và hủy & YC-?? & Đạt\\ \hline
TC-15 & Hỗ trợ các interface ERC-721 và AccessControl & YC-?? & Đạt\\ \hline
\multicolumn{4}{|l|}{\textit{Pool ERC-20 (giao dịch từng phần)}}\\ \hline
TC-16 & Nạp NFT vào pool nhận đúng số token, NFT bị khóa trong pool & YC-?? & Đạt\\ \hline
TC-17 & Nạp bằng safeTransferFrom trực tiếp; lô đã hủy một phần chỉ nhận phần còn lại & YC-?? & Đạt\\ \hline
TC-18 & Pool từ chối NFT không phải của hợp đồng tín chỉ & YC-?? & Đạt\\ \hline
TC-19 & Chuyển token ERC-20 từng phần cho người mua & YC-?? & Đạt\\ \hline
TC-20 & Người mua hủy từ pool, chứng nhận ghi người mua và dãy sê-ri; token đã đốt không dùng lại được (sử dụng kép) & YC-?? & Đạt\\ \hline
TC-21 & Bất biến: tổng cung token bằng tổng phần còn lại của các lô trong pool & YC-?? & Đạt\\ \hline
TC-22 & Đổi lại (redeem) lô cần đốt đủ token bằng phần còn lại & YC-?? & Đạt\\ \hline
TC-23 & Thao tác với lô không nằm trong pool bị từ chối & YC-?? & Đạt\\ \hline
TC-24 & NFT trong pool không thể bị chủ cũ bán riêng (bán kép) & YC-?? & Đạt\\ \hline
TC-25 & Hủy từ pool quá phần còn lại của lô bị từ chối, số dư không đổi & YC-?? & Đạt\\ \hline
TC-26 & Nạp lô đã hủy hết (không còn tồn tại) bị từ chối & YC-?? & Đạt\\ \hline
\end{longtable}
\endgroup}

\newcommand{\BangTestAPI}{%
\begingroup\small
\renewcommand{\arraystretch}{1.12}
\begin{longtable}{|p{1.35cm}|p{8.7cm}|p{1.55cm}|p{1.55cm}|}
\caption{Kết quả kiểm thử API và bộ đồng bộ sự kiện}\label{tab:testapi}\\
\hline
\textbf{Mã} & \textbf{Kịch bản kiểm thử} & \textbf{Yêu cầu} & \textbf{Kết quả}\\ \hline
\endfirsthead
\hline
\textbf{Mã} & \textbf{Kịch bản kiểm thử} & \textbf{Yêu cầu} & \textbf{Kết quả}\\ \hline
\endhead
API-01 & Đăng nhập bằng chữ ký ví, tài khoản mới tự chọn vai trò & YC-?? & Đạt\\ \hline
API-02 & Chữ ký sai hoặc dùng lại nonce bị từ chối & YC-?? & Đạt\\ \hline
API-03 & Không tự đăng ký được vai trò VERIFIER; ADMIN cấp được & YC-?? & Đạt\\ \hline
API-04 & Endpoint quản trị từ chối người không phải ADMIN và người chưa đăng nhập & YC-?? & Đạt\\ \hline
API-05 & Chủ dự án đăng ký dự án; trùng mã dự án bị từ chối & YC-?? & Đạt\\ \hline
API-06 & Gửi yêu cầu cấp tín chỉ cần báo cáo giám sát đã tải lên & YC-?? & Đạt\\ \hline
API-07 & Chỉ VVB duyệt được; duyệt xong trả về tham số mint với sê-ri đề xuất & YC-?? & Đạt\\ \hline
API-08 & Báo cáo thẩm định đã dùng không duyệt được yêu cầu khác (phát hành kép, off-chain) & YC-?? & Đạt\\ \hline
API-09 & Sự kiện CreditMinted tạo đợt cấp, khối tín chỉ và chuyển yêu cầu sang MINTED & YC-?? & Đạt\\ \hline
API-10 & Sự kiện trùng (cùng tx\_hash, log\_index) không bị ghi hai lần & YC-?? & Đạt\\ \hline
API-11 & Sự kiện CreditRetired tách khối và tạo chứng nhận tra cứu được & YC-?? & Đạt\\ \hline
API-12 & Tra cứu theo số sê-ri cho biết sê-ri đã hủy hay còn hiệu lực & YC-?? & Đạt\\ \hline
API-13 & Dashboard so sánh tín chỉ đã hủy với phần phát thải vượt hạn ngạch & YC-?? & Đạt\\ \hline
API-14 & Hủy hết phần còn lại thì không còn khối Held; thống kê khớp & YC-?? & Đạt\\ \hline
API-15 & Xuất dữ liệu dự án theo cấu trúc CAD Trust & YC-?? & Đạt\\ \hline
API-16 & Metadata NFT và tài liệu công khai cho phép đối chiếu mã băm & YC-?? & Đạt\\ \hline
E2E-01 & Bộ đồng bộ đọc sự kiện thật trên mạng Hardhat (mint, nạp pool, hủy); chạy lại không ghi trùng; phần còn lại on-chain khớp CSDL & YC-?? & Đạt\\ \hline
E2E-02 & Kiểm thử toàn hệ thống đang chạy: đăng nhập ví, đăng ký dự án, yêu cầu, VVB duyệt và mint, nạp pool, bán, hủy, tra cứu, dashboard & YC-?? & Đạt\\ \hline
\end{longtable}
\endgroup}

\newcommand{\BangDoPhu}{%
\begin{table}[htbp]
\centering
\caption{Độ phủ mã của bộ kiểm thử hợp đồng (solidity-coverage)}
\label{tab:dophu}
\begin{tabular}{|l|r|r|r|r|}
\hline
\textbf{Tệp} & \textbf{Câu lệnh} & \textbf{Nhánh} & \textbf{Hàm} & \textbf{Dòng}\\ \hline
CarbonCreditNFT.sol & 100\% & 94,44\% & 100\% & 100\%\\ \hline
CarbonPool.sol & 100\% & 75,00\% & 100\% & 100\%\\ \hline
\textbf{Tổng} & \textbf{100\%} & \textbf{88,46\%} & \textbf{100\%} & \textbf{100\%}\\ \hline
\end{tabular}
\end{table}}

\newcommand{\BangGas}{%
\begin{table}[htbp]
\centering
\caption{Chi phí gas của các thao tác chính (Solidity 0.8.24, optimizer 200 runs)}
\label{tab:gas}
\small
\begin{tabular}{|l|l|r|r|r|}
\hline
\textbf{Hợp đồng} & \textbf{Thao tác} & \textbf{Nhỏ nhất} & \textbf{Lớn nhất} & \textbf{Trung bình}\\ \hline
CarbonCreditNFT & Triển khai & -- & -- & 2.431.505\\ \hline
CarbonPool & Triển khai & -- & -- & 1.167.804\\ \hline
CarbonCreditNFT & mintBatch & 260.483 & 294.683 & 288.105\\ \hline
CarbonCreditNFT & retire (trực tiếp) & 206.215 & 261.001 & 243.049\\ \hline
CarbonCreditNFT & transferFrom & -- & -- & 57.310\\ \hline
CarbonCreditNFT & safeTransferFrom vào pool & -- & -- & 182.071\\ \hline
CarbonPool & deposit & 116.482 & 184.882 & 157.522\\ \hline
CarbonPool & transfer (PCC) & 29.683 & 51.595 & 47.208\\ \hline
CarbonPool & retire (từ pool) & 262.346 & 282.681 & 274.041\\ \hline
CarbonPool & redeem & -- & -- & 84.162\\ \hline
\end{tabular}
\end{table}}

\newcommand{\BangTinhToanKep}{%
\begin{table}[htbp]
\centering
\caption{Kịch bản chống tính toán kép và cơ chế chặn}
\label{tab:tinhtoankep}
\small
\renewcommand{\arraystretch}{1.2}
\begin{tabularx}{\textwidth}{|p{2.6cm}|X|X|p{1.9cm}|}
\hline
\textbf{Dạng} & \textbf{Kịch bản tấn công} & \textbf{Cơ chế chặn} & \textbf{Test}\\ \hline
Phát hành kép & VVB mint lô có dãy sê-ri chồng lên lô đã cấp của cùng dự án và vintage & nextAvailableSerial: sê-ri phải tăng dần, lỗi SerialOverlap & TC-03\\ \hline
Phát hành kép & Dùng lại một báo cáo thẩm định để cấp lần hai (kể cả cho dự án khác) & documentUsed on-chain (DocumentAlreadyUsed) và verifications.document\_hash unique off-chain & TC-05, API-08\\ \hline
Sử dụng kép & Hủy lại hoặc chuyển nhượng tín chỉ đã hủy & Trừ phần đã hủy khỏi lô; hủy hết thì đốt NFT & TC-09, TC-10\\ \hline
Sử dụng kép & Dùng lại token PCC đã đốt khi hủy từ pool & \_burn trước khi gọi retireFor; số dư không đủ thì giao dịch bị hoàn tác & TC-20\\ \hline
Bán kép & Chủ cũ bán riêng NFT đã nạp vào pool (đã nhận token tương ứng) & NFT thuộc sở hữu của pool, chỉ ra qua redeem khi đốt đủ token & TC-24, TC-21\\ \hline
Tuyên bố kép (trong hệ thống) & Hai bên cùng nhận đã bù trừ bằng một dãy sê-ri & Mỗi chứng nhận ghi dãy sê-ri và người thụ hưởng duy nhất; tra cứu công khai theo sê-ri & TC-08, API-12\\ \hline
\end{tabularx}
\end{table}}


% Tổng quan (Normal, không tiêu đề): chương trình bày các công nghệ đồ án dùng để hiện thực thiết kế ở
% Chương 3, kết quả cài đặt, kiểm thử và triển khai trên Sepolia.
% [VIẾT SAU] \canviet{Đoạn tổng quan chương}

% ============ PHẦN A — CÔNG NGHỆ SỬ DỤNG (từ Chương 2 cũ) ============
% [VIẾT SAU] Toàn bộ các mục dưới đây chưa viết nội dung — bỏ dấu % để khôi phục.
% % ------------------------------------------------------------
% \section{Blockchain và nền tảng Ethereum}
% \label{sec:blockchain}
% \label{sec:ethereum}
%
% \subsection*{Blockchain}
% \label{subsec:cautrucblock}
% \label{subsec:dongthuan}
% % Ý cần có (ngắn gọn, chỉ phần đồ án dùng): khối, mã băm khối trước, sổ cái chỉ ghi thêm (ledger) ->
% % vì sao dữ liệu on-chain khó sửa, phục vụ yêu cầu toàn vẹn ở mục \ref{sec:phichucnang} và VĐ-2;
% % đồng thuận PoW/PoS, Ethereum chuyển sang PoS (The Merge 15/9/2022).
% % Trích: \cite{nakamoto2008bitcoin,zheng2017overview,ethereum2024pos}. Nhắc \autoref{fig:chuoikhoi} trước hình.
% % [VIẾT SAU] \canviet{Blockchain}
% %
% % \HinhChuoiKhoi
%
% \subsection*{Lựa chọn nền tảng blockchain}
% \label{subsec:phanloai}
% % Bảng so sánh Ethereum (public) với Hyperledger Fabric (permissioned, mà Vaccargiu dùng) và Chia
% % (CAD Trust dùng): mô hình truy cập, đồng thuận, ngôn ngữ smart contract, chuẩn token, ví, testnet, chi phí.
% % Lý do chọn Ethereum: phiếu giao nhiệm vụ; chuẩn ERC-721/ERC-20; MetaMask; testnet công khai để bên thứ ba
% % tự kiểm tra (VĐ-2, NC-07). Nhược điểm: phí gas, dữ liệu công khai -> xử lý ở yêu cầu quyền riêng tư.
% % [VIẾT SAU] \canviet{Lựa chọn nền tảng blockchain}
%
% \subsection*{Tài khoản, giao dịch và phí gas}
% \label{subsec:taikhoan}
% \label{subsec:evm}
% % Ý cần có: EOA và contract account; khóa riêng, địa chỉ ví; giao dịch ký số, nonce; EVM, gas, EIP-1559.
% % Đây là cơ sở cho phần đo gas ở Chương 3. Trích: \cite{buterin2014ethereum,wood2014yellow,eip1559}.
% % [VIẾT SAU] \canviet{Tài khoản, giao dịch và phí gas}
%
% \subsection*{Testnet Sepolia và Etherscan}
% \label{subsec:sepolia}
% % Ý cần có: mainnet và testnet; Sepolia, faucet; Etherscan để tra giao dịch, xác minh (verify) mã hợp đồng.
% % Lý do: phạm vi nguyên mẫu (mục \ref{sec:muctieu}), không tốn phí thật.
% % [VIẾT SAU] \canviet{Testnet Sepolia}
%
% % ------------------------------------------------------------
% \section{Smart contract và công cụ phát triển}
% \label{sec:smartcontract}
%
% \subsection*{Smart contract và ngôn ngữ Solidity}
% \label{subsec:khainiemsc}
% \label{subsec:solidity}
% % Ý cần có: định nghĩa (Szabo); tự thực thi, không sửa được mã sau triển khai; event và log (dùng cho
% % bộ đồng bộ sự kiện); cấu trúc contract (state variable, function, modifier, event), kế thừa.
% % Yêu cầu phục vụ: NC-03, NC-06 (ràng buộc sê-ri, không hủy hai lần thành ràng buộc kỹ thuật).
% % Trích: \cite{szabo1997formalizing,buterin2014ethereum,solidity2024docs}.
% % [VIẾT SAU] \canviet{Smart contract và Solidity}
%
% \subsection*{Thư viện OpenZeppelin và Hardhat}
% \label{subsec:openzeppelin}
% % OpenZeppelin 5: ERC721, ERC20, AccessControl (VERIFIER_ROLE...), Pausable, ReentrancyGuard -> NC-09 và
% % yêu cầu bảo mật. Hardhat: biên dịch, Hardhat Network, test Mocha/Chai, coverage, gas reporter.
% % Phương án thay thế: Foundry, Truffle (đã ngừng phát triển); lý do chọn Hardhat (JavaScript cùng ngôn ngữ backend).
% % Trích: \cite{openzeppelin2024contracts}.
% % [VIẾT SAU] \canviet{OpenZeppelin và Hardhat}
%
% \subsection*{Các lỗ hổng thường gặp của smart contract}
% \label{subsec:baomat}
% % Reentrancy, kiểm soát truy cập sai, tràn số (Solidity >= 0.8 tự kiểm tra), tx.origin. Mỗi lỗ hổng
% % 1–2 câu và cách phòng trong đồ án; công cụ quét Slither. Trích: \cite{atzei2017survey}.
% % [VIẾT SAU] \canviet{Lỗ hổng smart contract}
%
% % ------------------------------------------------------------
% \section{Các chuẩn token}
% \label{sec:token-standards}
%
% \subsection*{ERC-20 và ERC-721}
% \label{subsec:erc20}
% \label{subsec:erc721}
% % ERC-20: fungible token, totalSupply, balanceOf, transfer, approve, transferFrom, allowance \cite{eip20}.
% % ERC-721: NFT, ownerOf, safeTransferFrom, tokenURI, onERC721Received \cite{eip721}.
% % [VIẾT SAU] \canviet{ERC-20 và ERC-721}
%
% \subsection*{Lựa chọn chuẩn token cho tín chỉ các-bon}
% \label{subsec:chontoken}
% % Bảng so sánh ERC-20 / ERC-721 / ERC-1155 theo: giữ metadata riêng từng lô, chia nhỏ, thanh khoản,
% % hỗ trợ ví. Kết luận: mỗi lô là một NFT ERC-721 (giữ dự án, vintage, dãy sê-ri, NC-03), pool ERC-20
% % cho giao dịch từng phần (NC-05); kế thừa mô hình Batch NFT/TCO2/BCT của Toucan (mục \ref{sec:khaosatsobo}).
% % [VIẾT SAU] \canviet{Lựa chọn chuẩn token}
%
% % ------------------------------------------------------------
% \section{Ứng dụng web giao tiếp với blockchain}
% \label{sec:dapp}
%
% \subsection*{Kiến trúc dApp}
% \label{subsec:kientrucdapp}
% % Luồng một giao dịch từ giao diện tới smart contract (giải thích \autoref{fig:luongdapp}).
% % [VIẾT SAU] \canviet{Kiến trúc dApp}
% %
% % \HinhLuongDApp
%
% \subsection*{Ví MetaMask, ethers.js và Next.js}
% \label{subsec:metamask}
% \label{subsec:ethers}
% % MetaMask giữ khóa riêng, EIP-1193 (window.ethereum) \cite{eip1193}; ethers.js: Provider, Signer,
% % Contract, ABI, call và send \cite{ethers2024docs}; Next.js cho giao diện. Yêu cầu phục vụ: tính dễ dùng,
% % người dùng tự ký giao dịch (backend không giữ khóa). Phương án thay thế: web3.js, viem; React thuần.
% % [VIẾT SAU] \canviet{MetaMask, ethers.js và Next.js}
%
% % ------------------------------------------------------------
% \section{Công nghệ phía off-chain}
% \label{sec:offchain-congnghe}
% \label{sec:onoffchain-lythuyet}
% % Node.js + Express (API RESTful), PostgreSQL (dữ liệu theo CAD Trust V2.0, NC-10), JWT, xác thực bằng
% % chữ ký ví. Tiêu chí chọn dữ liệu on-chain/off-chain theo \cite{xu2017taxonomy}; neo mã băm keccak256 của hồ sơ;
% % IPFS \cite{benet2014ipfs} là hướng mở rộng (hiện lưu tệp cục bộ). Phương án thay thế: MongoDB; lý do chọn
% % CSDL quan hệ (mô hình CAD Trust có quan hệ rõ, giống UNDP dùng PostgreSQL).
% % [VIẾT SAU] \canviet{Công nghệ off-chain}


% ============ PHẦN B — XÂY DỰNG, KIỂM THỬ, TRIỂN KHAI (từ Chương 3 cũ) ============
% ------------------------------------------------------------
\section{Xây dựng ứng dụng}
\label{sec:xaydung}

\subsection*{Thư viện và công cụ sử dụng}
\label{sec:moitruong}
% Mẫu SoICT: bảng công cụ có phiên bản (và URL). Giải thích \autoref{tab:congnghe}; lý do chọn công cụ đã nêu
% ở Chương 2 nên chỉ nhắc lại ngắn. Cấu trúc mã nguồn theo \autoref{fig:cautrucma}.
% [VIẾT SAU] \canviet{Thư viện và công cụ sử dụng}

\BangCongNghe

\HinhCauTrucMa

\subsection*{Kết quả đạt được}
\label{subsec:ketquadatduoc}
\label{sec:caidat-hopdong}
\label{sec:caidat-backend}
% Mẫu SoICT: mô tả các sản phẩm đóng gói, thành phần, vai trò; bảng thống kê số dòng mã, số tệp, số hợp đồng,
% số endpoint, dung lượng.
% Nội dung cài đặt từ Chương 4 cũ (chỉ tóm lược, chi tiết đóng góp ở Chương 4):
%  - CarbonCreditNFT: struct CreditBatch (projectId, vintage, serialFrom, quantity, retired, documentHash,
%    verifier, issuedAt); mintBatch: seriesKey = keccak256(projectId, vintage), serialFrom >= nextAvailableSerial,
%    documentUsed[documentHash]; retire hủy từ đầu dãy sê-ri còn lại; retireFor cho pool (RETIREMENT_AGENT_ROLE);
%    Pausable ghi đè _update.
%  - CarbonPool: decimals = 0; nạp qua onERC721Received; redeem đốt remainingOf(tokenId); retire đốt token rồi
%    gọi retireFor; ReentrancyGuard.
%  - Phân quyền: DEFAULT_ADMIN_ROLE và vai trò phát hành ISSUER_ROLE (cơ quan quản lý, mỗi tỉnh một tài khoản),
%    VERIFIER_ROLE (VVB, chỉ phê duyệt off-chain), RETIREMENT_AGENT_ROLE (pool). [CẦN SỬA MÃ] thêm ISSUER_ROLE.
%  - Backend: db/schema.sql từ ERD; xác thực nonce -> personal_sign -> ethers.verifyMessage -> JWT 8 giờ;
%    API theo \autoref{tab:api} (backend không giữ khóa riêng); indexer quét queryFilter, khóa (tx_hash, log_index),
%    chưa xử lý reorg.
% [VIẾT SAU] \canviet{Kết quả đạt được}

% \BangAPI   % [17/11/2026] đã dùng ở Chương 3, mục Thiết kế API

\subsection*{Minh họa các chức năng chính}
\label{subsec:minhhoa}
\label{sec:caidat-frontend}
% Mỗi màn hình 1 ảnh chụp + giải thích ngắn: danh sách dự án, chi tiết dự án, tra cứu sê-ri, chứng nhận hủy,
% dự án của tôi, thẩm định (VVB), ví tín chỉ, dashboard, quản trị. Chụp khi chạy thật (README "Chạy thử cục bộ").
% [VIẾT SAU] \canviet{Minh họa các chức năng chính}

% ------------------------------------------------------------
\section{Kiểm thử}
\label{sec:kiemthu}
\label{sec:pp-kiemthu}
\label{sec:kq-chucnang}
% Mẫu SoICT: 2–3 trang; thiết kế ca kiểm thử cho 2–3 chức năng quan trọng nhất (phát hành, hủy từ pool, tra cứu),
% nêu kỹ thuật kiểm thử; tổng kết số ca và kết quả; giải thích ca không đạt. Ca kiểm thử khác đưa vào Phụ lục.
% Ba mức: kiểm thử đơn vị hợp đồng trên Hardhat Network (Mocha/Chai, loadFixture); kiểm thử API với pg-mem và
% supertest; kiểm thử tích hợp e2e-indexer và smoke-api. Ghi cấu hình máy, phiên bản Node, Hardhat.
% Nhận xét độ phủ: nhánh của CarbonPool 75% vì ... (xem coverage/index.html).
% Chi phí gas: quy đổi ETH/USD theo giá gas tại thời điểm đo (ghi nguồn, ngày); vì sao mintBatch, retire tốn gas.
% Kịch bản chống tính toán kép (\BangTinhToanKep) đặt ở Chương 4, mục \ref{sec:gp-tinhtoankep}.
% [VIẾT SAU] \canviet{Kiểm thử}

\BangTestHopDong

\BangTestAPI

\BangDoPhu

\BangGas

% ------------------------------------------------------------
\section{Triển khai}
\label{sec:trienkhai}
\label{sec:kq-sepolia}
% Giải thích \autoref{fig:trienkhai}: Sepolia + RPC, backend + PostgreSQL, frontend; địa chỉ hợp đồng,
% link Etherscan (đã verify), vài giao dịch mẫu (mint, retire), ảnh chụp.
% [VIẾT SAU] \canviet{Triển khai}

\HinhTrienKhai


% ------------------------------------------------------------
% Kết chương (Normal, không tiêu đề): tóm tắt bộ công nghệ (bảng Công nghệ | Phiên bản | Vai trò | Yêu cầu
% phục vụ), kết quả cài đặt, kiểm thử, triển khai; dẫn sang Chương 5 (giải pháp và đóng góp).
% [VIẾT SAU] \canviet{Kết chương 4}
